\documentclass[11pt,a4paper,reqno]{amsart}
\usepackage[T1]{fontenc}
\usepackage{lmodern}
\usepackage[margin=25mm]{geometry}
\usepackage{amsmath,listings}
\usepackage[hidelinks]{hyperref}
\hypersetup{pdftitle={Supplementary tables: 131 six-color partitions of [1,537]},pdfauthor={Zhiyang Chen, Ligeng Zhu, Fan Long}}
\lstset{basicstyle=\scriptsize\ttfamily,columns=fullflexible,keepspaces=true,
  breaklines=true,showstringspaces=false,aboveskip=1pt,belowskip=5pt}
\title[Supplementary partitions]{Supplementary Tables: 131 Six-Color Sum-Free Partitions of $[1,537]$}
\author{Zhiyang Chen}
\address{Department of Computer Science, University of Toronto, Toronto, Canada}
\email{zhiychen@cs.toronto.edu}
\author{Ligeng Zhu}
\address{NVIDIA Research}
\email{ligengz@nvidia.com}
\author{Fan Long}
\address{Department of Computer Science, University of Toronto, Toronto, Canada}
\email{fanl@cs.toronto.edu}
\date{22 October 2026}
\begin{document}
\maketitle
Each word gives $c(1)c(2)\cdots c(537)$, with colors numbered 1--6 in the order
of their least elements. Spaces and line breaks have no mathematical meaning.
These are the 131 partitions described in the accompanying article.
Their classes are sum-free modulo $538$, and the partitions are pairwise
inequivalent under color renaming and multiplication by units modulo $538$.
\par\medskip
\subsection*{Family $\mathcal F_1$ (107 partitions)}

\noindent\begin{minipage}{\linewidth}\textbf{F1.001}
\begin{lstlisting}
1234122334 4143142213 1213132413 3144432514 2215123515 2535454535 5535325532 5555535545 2355535322
5344145342 2153215122 4122344413 3142213121 3122413414 4332214321 6123412233 4414314221 3151613641
3366463356 6236655666 5643666656 6566666666 6666666666 6666656666 6666666666 6666666666 6666666666
6666666666 6666666666 6666666666 6566666666 6666666666 6666656656 6663465666 5566326653 3646633146
3161513122 4134144332 2143216123 4122334414 3142213121 3122413314 4432214221 5123512243 5414435223
5355532545 5355555235 5235355535 4545352515 3215122415 2344413314 2313121312 2413414433 2214321
\end{lstlisting}
\end{minipage}\par\smallskip

\noindent\begin{minipage}{\linewidth}\textbf{F1.002}
\begin{lstlisting}
1234122334 4143142213 1213122413 3144433514 2215123515 2535454535 5535225532 5555535555 2355535322
5344145342 2153215122 4122344413 3142213121 3122413414 4332214321 6123412233 4414314221 3151613641
3366463356 6256655666 5653666656 6566666666 6666666666 6666656666 6666666666 6666666666 6666666666
6666666666 6666666666 6666666666 6566666666 6666666666 6666656656 6663565666 5566526653 3646633146
3161513122 4134144332 2143216123 4122334414 3142213121 3122413314 4432214221 5123512243 5414435223
5355532555 5355555235 5225355535 4545352515 3215122415 3344413314 2213121312 2413414433 2214321
\end{lstlisting}
\end{minipage}\par\smallskip

\noindent\begin{minipage}{\linewidth}\textbf{F1.003}
\begin{lstlisting}
1234122334 4143142213 1213132413 3144432514 2215123515 2435454535 5535325532 5555535555 5355535322
5344145342 2153215122 4122344413 3142213121 3122413414 4332214221 6123412233 4414314221 3151613641
3366463356 6236655666 5653666656 6566666666 6666646666 6666656666 6666666666 6666666666 6666666666
6666666666 6666666666 6666666666 6566666666 6466666666 6666656656 6663565666 5566326653 3646633146
3161513122 4134144332 2143216122 4122334414 3142213121 3122413314 4432214221 5123512243 5414435223
5355535555 5355555235 5235355535 4545342515 3215122415 2344413314 2313121312 2413414433 2214321
\end{lstlisting}
\end{minipage}\par\smallskip

\noindent\begin{minipage}{\linewidth}\textbf{F1.004}
\begin{lstlisting}
1234122334 4143142213 1213132413 3144432514 2215123515 2535454535 5535225532 5555535555 5355535322
5344145342 2153215122 4122344413 3142213121 3122413414 4332214221 6123412233 4414314221 3151613641
3366463356 6256655666 5653666656 6566666666 6666666666 6666656666 6666666666 6666666666 6666666666
6666666666 6666666666 6666666666 6566666666 6666666666 6666656656 6663565666 5566526653 3646633146
3161513122 4134144332 2143216122 4122334414 3142213121 3122413314 4432214221 5123512243 5414435223
5355535555 5355555235 5225355535 4545352515 3215122415 2344413314 2313121312 2413414433 2214321
\end{lstlisting}
\end{minipage}\par\smallskip

\noindent\begin{minipage}{\linewidth}\textbf{F1.005}
\begin{lstlisting}
1234122334 4143142213 1213122413 3144433514 2215123515 2435454535 5535225532 5455335555 2355535322
5344145342 2153215122 4122344413 3142213121 3122413414 4332214321 6123412233 4414314221 3151613641
3366463356 6236655666 5643666656 6566666666 6666666666 6666656666 6666666666 6666666666 6666666666
6666666666 6666666666 6666666666 6566666666 6666666666 6666656656 6663465666 5566326653 3646633146
3161513122 4134144332 2143216123 4122334414 3142213121 3122413314 4432214221 5123512243 5414435223
5355532555 5335545235 5225355535 4545342515 3215122415 3344413314 2213121312 2413414433 2214321
\end{lstlisting}
\end{minipage}\par\smallskip

\noindent\begin{minipage}{\linewidth}\textbf{F1.006}
\begin{lstlisting}
1234122334 4143142213 1213132413 3144432514 2215123515 2535454535 5535225532 5555535555 5355535322
5344145342 2153215122 4122344413 3142213121 3122413414 4332214321 6123412233 4414314221 6151613641
3364463356 6236655666 5663666656 6566636666 6666666666 6666656666 6666666666 6666666666 6666666666
6666666666 6666666666 6666666666 6566666666 6666666666 6366656656 6663665666 5566326653 3644633146
3161516122 4134144332 2143216123 4122334414 3142213121 3122413314 4432214221 5123512243 5414435223
5355535555 5355555235 5225355535 4545352515 3215122415 2344413314 2313121312 2413414433 2214321
\end{lstlisting}
\end{minipage}\par\smallskip

\noindent\begin{minipage}{\linewidth}\textbf{F1.007}
\begin{lstlisting}
1234122334 4143142213 1213122413 3144432514 2215123515 2535454535 2535225532 5555335545 2355234322
5344145342 2153215122 4122344413 3142213121 3122413414 4332214321 6123412233 4414314221 3151613641
3366463356 6236655666 5553666656 6566666666 6666666666 6666656666 6666666666 6666666666 6666666666
6666666666 6666666666 6666666666 6566666666 6666666666 6666656656 6663555666 5566326653 3646633146
3161513122 4134144332 2143216123 4122334414 3142213121 3122413314 4432214221 5123512243 5414435223
4325532545 5335555235 5225352535 4545352515 3215122415 2344413314 2213121312 2413414433 2214321
\end{lstlisting}
\end{minipage}\par\smallskip

\noindent\begin{minipage}{\linewidth}\textbf{F1.008}
\begin{lstlisting}
1234122334 4143142213 1213122413 3144432514 2215123515 2535454535 2535325532 5555335545 5355535322
5344145342 2153215122 4122344413 3142213121 3122413414 4332214321 6123412233 4414314221 3151613661
3364463356 6236655666 5663666556 6566666666 6666666666 6666656666 6666666666 6666666666 6666666666
6666666666 6666666666 6666666666 6566666666 6666666666 6666656655 6663665666 5566326653 3644633166
3161513122 4134144332 2143216123 4122334414 3142213121 3122413314 4432214221 5123512243 5414435223
5355535545 5335555235 5235352535 4545352515 3215122415 2344413314 2213121312 2413414433 2214321
\end{lstlisting}
\end{minipage}\par\smallskip

\noindent\begin{minipage}{\linewidth}\textbf{F1.009}
\begin{lstlisting}
1234122334 4143142213 1213122413 3144432514 2215123515 2535454535 5535225533 5555535555 5355535322
5344145342 2153215122 4122344413 3142213121 3122413414 4332214221 6123412233 4414314221 3151613641
3366463356 6256655666 5653666656 6566666666 6666666666 6666656666 6666666666 6666666666 6666666666
6666666666 6666666666 6666666666 6566666666 6666666666 6666656656 6663565666 5566526653 3646633146
3161513122 4134144332 2143216122 4122334414 3142213121 3122413314 4432214221 5123512243 5414435223
5355535555 5355555335 5225355535 4545352515 3215122415 2344413314 2213121312 2413414433 2214321
\end{lstlisting}
\end{minipage}\par\smallskip

\noindent\begin{minipage}{\linewidth}\textbf{F1.010}
\begin{lstlisting}
1234122334 4143142213 1213132413 3144432514 2215123515 2435454535 5535225533 5555335555 5355535322
5344145342 2153215122 4122344413 3142213121 3122413414 4332214321 6123412233 4414314221 3151613641
3366463356 6256655666 5653666656 6566666666 6666666666 6666656666 6666666666 6666666666 6666666666
6666666666 6666666666 6666666666 6566666666 6666666666 6666656656 6663565666 5566526653 3646633146
3161513122 4134144332 2143216123 4122334414 3142213121 3122413314 4432214221 5123512243 5414435223
5355535555 5335555335 5225355535 4545342515 3215122415 2344413314 2313121312 2413414433 2214321
\end{lstlisting}
\end{minipage}\par\smallskip

\noindent\begin{minipage}{\linewidth}\textbf{F1.011}
\begin{lstlisting}
1234122334 4143142213 1213132413 3144432514 2315123515 2535454535 5535225532 5555335555 5355535322
5344145342 2153215132 4122344413 3142213121 3122413414 4332214321 6123412233 4414314221 3151613641
3366463356 6256655666 5653666656 6566666666 6666666666 6666656666 6666666666 6666666666 6666666666
6666666666 6666666666 6666666666 6566666666 6666666666 6666656656 6663565666 5566526653 3646633146
3161513122 4134144332 2143216123 4122334414 3142213121 3122413314 4432214231 5123512243 5414435223
5355535555 5335555235 5225355535 4545352515 3215132415 2344413314 2313121312 2413414433 2214321
\end{lstlisting}
\end{minipage}\par\smallskip

\noindent\begin{minipage}{\linewidth}\textbf{F1.012}
\begin{lstlisting}
1234122334 4143142213 1213122413 3144432514 2215123515 2435454535 5535325532 5455535545 5355535322
5344145342 2153215132 4122344413 3142213121 3122413414 4332214321 6123412233 4414314221 6151613641
3364463356 6236655666 5663666656 6566656666 6666666666 6666656666 6666666666 6666666666 6666666666
6666666666 6666666666 6666666666 6566666666 6666666666 6566656656 6663665666 5566326653 3644633146
3161516122 4134144332 2143216123 4122334414 3142213121 3122413314 4432214231 5123512243 5414435223
5355535545 5355545235 5235355535 4545342515 3215122415 2344413314 2213121312 2413414433 2214321
\end{lstlisting}
\end{minipage}\par\smallskip

\noindent\begin{minipage}{\linewidth}\textbf{F1.013}
\begin{lstlisting}
1234122334 4143142213 1213122413 3144433514 2215123515 2435454535 5535225532 5555535555 3355535322
5344145342 2153215122 4122344413 3142213121 3122413414 4332214221 6123412233 4414316221 3151613641
3366463356 6236655666 5643666656 6566666665 6666666666 6666656666 6666666666 6666666666 6666666666
6666666666 6666666666 6666666666 6566666666 6666666566 6666656656 6663465666 5566326653 3646633146
3161513122 6134144332 2143216122 4122334414 3142213121 3122413314 4432214221 5123512243 5414435223
5355533555 5355555235 5225355535 4545342515 3215122415 3344413314 2213121312 2413414433 2214321
\end{lstlisting}
\end{minipage}\par\smallskip

\noindent\begin{minipage}{\linewidth}\textbf{F1.014}
\begin{lstlisting}
1234122334 4143142213 1213122413 3144433514 2215123515 2535454535 2535325532 5555535555 5355534322
5344145342 2153215122 4122344413 3142213121 3122413414 4332214321 6123412233 4414316221 3151613641
3366463356 6236655666 5653666656 6566666665 6666666666 6666656666 6666666666 6666666666 6666666666
6666666666 6666666666 6666666666 6566666666 6666666566 6666656656 6663565666 5566326653 3646633146
3161513122 6134144332 2143216123 4122334414 3142213121 3122413314 4432214221 5123512243 5414435223
4355535555 5355555235 5235352535 4545352515 3215122415 3344413314 2213121312 2413414433 2214321
\end{lstlisting}
\end{minipage}\par\smallskip

\noindent\begin{minipage}{\linewidth}\textbf{F1.015}
\begin{lstlisting}
1234122334 4143142213 1213132413 3144432514 2215123515 2435454535 2535325532 5455535555 5355535322
5344145342 2153215122 4122344413 3142213121 3122413414 4332214321 6123412233 4414314221 6151613641
3364463356 6236655666 5663666656 6566636666 6666666666 6666656666 6666666666 6666666666 6666666666
6666666666 6666666666 6666666666 6566666666 6666666666 6366656656 6663665666 5566326653 3644633146
3161516122 4134144332 2143216123 4122334414 3142213121 3122413314 4432214221 5123512243 5414435223
5355535555 5355545235 5235352535 4545342515 3215122415 2344413314 2313121312 2413414433 2214321
\end{lstlisting}
\end{minipage}\par\smallskip

\noindent\begin{minipage}{\linewidth}\textbf{F1.016}
\begin{lstlisting}
1234122334 4143142213 1213132413 3144432514 2215123515 2535454535 5535225532 5555535545 5355535322
5344145342 2153215122 4122344413 3142213121 3122413414 4332214321 6123412233 4414314221 6151613661
3364463356 6236655666 5643665556 6566656666 6646666666 6666656666 6666666666 6666666666 6666666666
6666666666 6666666666 6666666666 6566666666 6666466666 6566656655 5663465666 5566326653 3644633166
3161516122 4134144332 2143216123 4122334414 3142213121 3122413314 4432214221 5123512243 5414435223
5355535545 5355555235 5225355535 4545352515 3215122415 2344413314 2313121312 2413414433 2214321
\end{lstlisting}
\end{minipage}\par\smallskip

\noindent\begin{minipage}{\linewidth}\textbf{F1.017}
\begin{lstlisting}
1234122334 4143142213 1213132413 3144432514 2215123515 2535454535 5535325532 5555535555 2355535322
5344145342 2153215132 4122344413 3142213121 3122413414 4332214221 6123412233 4414314221 6151613661
3364463356 6236655666 5663666556 6566636666 6666666666 6666656666 6666666666 6666666666 6666666666
6666666666 6666666666 6666666666 6566666666 6666666666 6366656655 6663665666 5566326653 3644633166
3161516122 4134144332 2143216122 4122334414 3142213121 3122413314 4432214231 5123512243 5414435223
5355532555 5355555235 5235355535 4545352515 3215122415 2344413314 2313121312 2413414433 2214321
\end{lstlisting}
\end{minipage}\par\smallskip

\noindent\begin{minipage}{\linewidth}\textbf{F1.018}
\begin{lstlisting}
1234122334 4143142213 1213132413 3144432514 2315123515 2535454535 5535225532 5555535555 5355535322
5344145342 2153215122 4122344413 3142213121 3122413414 4332214221 6123412233 4414316221 3151613641
3366463356 6256655666 5653666656 6566666665 6666666666 6666656666 6666666666 6666666666 6666666666
6666666666 6666666666 6666666666 6566666666 6666666566 6666656656 6663565666 5566526653 3646633146
3161513122 6134144332 2143216122 4122334414 3142213121 3122413314 4432214221 5123512243 5414435223
5355535555 5355555235 5225355535 4545352515 3215132415 2344413314 2313121312 2413414433 2214321
\end{lstlisting}
\end{minipage}\par\smallskip

\noindent\begin{minipage}{\linewidth}\textbf{F1.019}
\begin{lstlisting}
1234122334 4143142213 1213122413 3144432514 2215123515 2435454535 5535225533 5555535555 5355535322
5344145342 2153215122 4122344413 3142213121 3122413414 4332214321 6123412233 4414314221 3151613661
3364463356 6236655666 5663666556 6566666666 6666666666 6666656666 6666666666 6666666666 6666666666
6666666666 6666666666 6666666666 6566666666 6666666666 6666656655 6663665666 5566326653 3644633166
3161513122 4134144332 2143216123 4122334414 3142213121 3122413314 4432214221 5123512243 5414435223
5355535555 5355555335 5225355535 4545342515 3215122415 2344413314 2213121312 2413414433 2214321
\end{lstlisting}
\end{minipage}\par\smallskip

\noindent\begin{minipage}{\linewidth}\textbf{F1.020}
\begin{lstlisting}
1234122334 4143142213 1213122413 3144432514 2215123515 2535454535 5535225532 5555335555 2355235322
5344145342 2153215122 4122344413 3142213121 3122413414 4332214321 6123412233 4414314221 3151613641
3366463356 6256655666 5643665656 6566666665 6666666666 6336656666 6666666666 6666666666 6666666666
6666666666 6666666666 6666666666 6566336666 6666666566 6666656656 5663465666 5566526653 3646633146
3161513122 4134144332 2143216123 4122334414 3142213121 3122413314 4432214221 5123512243 5414435223
5325532555 5335555235 5225355535 4545352515 3215122415 2344413314 2213121312 2413414433 2214321
\end{lstlisting}
\end{minipage}\par\smallskip

\noindent\begin{minipage}{\linewidth}\textbf{F1.021}
\begin{lstlisting}
1234122334 4143142213 1213122413 3144432514 2215123515 2535454535 5535325532 5555535555 5355235322
5344145342 2153215122 4122344413 3142213121 3122413414 4332214321 6123412233 4414314221 6151613661
3364463356 6236655666 5663666556 6566656666 6666666666 6666656666 6666666666 6666666666 6666666666
6666666666 6666666666 6666666666 6566666666 6666666666 6566656655 6663665666 5566326653 3644633166
3161516122 4134144332 2143216123 4122334414 3142213121 3122413314 4432214221 5123512243 5414435223
5325535555 5355555235 5235355535 4545352515 3215122415 2344413314 2213121312 2413414433 2214321
\end{lstlisting}
\end{minipage}\par\smallskip

\noindent\begin{minipage}{\linewidth}\textbf{F1.022}
\begin{lstlisting}
1234122334 4143142213 1213122413 3144432514 2315123515 2535454535 5535225532 5455535555 2355535322
5344145342 2153215132 4122344413 3142213121 3122413414 4332214221 6123412233 4414314221 3151613651
3366463356 6256655666 5653666656 6566666666 6666666666 6666656666 6666666666 6666666666 6666666666
6666666666 6666666666 6666666666 6566666666 6666666666 6666656656 6663565666 5566526653 3646633156
3161513122 4134144332 2143216122 4122334414 3142213121 3122413314 4432214231 5123512243 5414435223
5355532555 5355545235 5225355535 4545352515 3215132415 2344413314 2213121312 2413414433 2214321
\end{lstlisting}
\end{minipage}\par\smallskip

\noindent\begin{minipage}{\linewidth}\textbf{F1.023}
\begin{lstlisting}
1234122334 4143142213 1213122413 3144432514 2315123515 2535454535 5535225532 5555535555 5355235322
5344145342 2153215122 4122344413 3142213121 3122413414 4332214221 6123412233 4414316221 3151613641
3366463356 6256655666 5643666656 6566666665 6666666666 6666656666 6666666666 6666666666 6666666666
6666666666 6666666666 6666666666 6566666666 6666666566 6666656656 6663465666 5566526653 3646633146
3161513122 6134144332 2143216122 4122334414 3142213121 3122413314 4432214221 5123512243 5414435223
5325535555 5355555235 5225355535 4545352515 3215132415 2344413314 2213121312 2413414433 2214321
\end{lstlisting}
\end{minipage}\par\smallskip

\noindent\begin{minipage}{\linewidth}\textbf{F1.024}
\begin{lstlisting}
1234122334 4143142213 1213122413 3144433514 2215123515 2435454535 5535225532 5555335555 5355535322
5344145342 2153215122 4122344413 3142213121 3122413414 4332214321 6123412233 4414314221 3151613641
2266463356 6236655666 5643665656 6566666666 6666666666 6666656666 6666666666 6666666666 6666666666
6666666666 6666666666 6666666666 6566666666 6666666666 6666656656 5663465666 5566326653 3646622146
3161513122 4134144332 2143216123 4122334414 3142213121 3122413314 4432214221 5123512243 5414435223
5355535555 5335555235 5225355535 4545342515 3215122415 3344413314 2213121312 2413414433 2214321
\end{lstlisting}
\end{minipage}\par\smallskip

\noindent\begin{minipage}{\linewidth}\textbf{F1.025}
\begin{lstlisting}
1234122334 4143142213 1213122413 3144433514 2215123515 2435454535 5535325532 5555535555 5355235322
5344145342 2153215132 4122344413 3142213121 3122413414 4332214321 6123412233 4414314221 3151613641
3366463356 6256655666 5553666656 6566666666 6666666666 6666656666 6666666666 6666666666 6666666666
6666666666 6666666666 6666666666 6566666666 6666666666 6666656656 6663555666 5566526653 3646633146
3161513122 4134144332 2143216123 4122334414 3142213121 3122413314 4432214231 5123512243 5414435223
5325535555 5355555235 5235355535 4545342515 3215122415 3344413314 2213121312 2413414433 2214321
\end{lstlisting}
\end{minipage}\par\smallskip

\noindent\begin{minipage}{\linewidth}\textbf{F1.026}
\begin{lstlisting}
1234122334 4143142213 1213122413 3144433514 2315123515 2535454535 5535325533 5555535555 5355534322
5344145342 2153215122 4122344413 3142213121 3122413414 4332214321 6123412233 4414314221 3151613641
3266463356 6256655666 5653666656 6566666666 6666666666 6666656666 6666666666 6666666666 6666666666
6666666666 6666666666 6666666666 6566666666 6666666666 6666656656 6663565666 5566526653 3646623146
3161513122 4134144332 2143216123 4122334414 3142213121 3122413314 4432214221 5123512243 5414435223
4355535555 5355555335 5235355535 4545352515 3215132415 3344413314 2213121312 2413414433 2214321
\end{lstlisting}
\end{minipage}\par\smallskip

\noindent\begin{minipage}{\linewidth}\textbf{F1.027}
\begin{lstlisting}
1234122334 4143142213 1213132413 3144432514 2215123515 2535454535 5535225532 5455335545 5355535322
5344145342 2153215132 4122344413 3142213121 3122413414 4332214321 6123412233 4414316221 3151613651
3366463356 6236655666 5653666656 6566666665 6664666666 6666656666 6666666666 6666666666 6666666666
6666666666 6666666666 6666666666 6566666666 6664666566 6666656656 6663565666 5566326653 3646633156
3161513122 6134144332 2143216123 4122334414 3142213121 3122413314 4432214231 5123512243 5414435223
5355535545 5335545235 5225355535 4545352515 3215122415 2344413314 2313121312 2413414433 2214321
\end{lstlisting}
\end{minipage}\par\smallskip

\noindent\begin{minipage}{\linewidth}\textbf{F1.028}
\begin{lstlisting}
1234122334 4143142213 1213132413 3144432514 2215123515 2535454535 5535225532 5555535555 5355535322
5344145342 2153215122 4122344413 3142213121 3122413414 4332214221 6123412233 4414314221 6151613641
3364463356 6236655666 5663666653 6566656666 6666663666 6666656666 6666666666 6666666666 6666666666
6666666666 6666666666 6666666666 6566666666 3666666666 6566656356 6663665666 5566326653 3644633146
3161516122 4134144332 2143216122 4122334414 3142213121 3122413314 4432214221 5123512243 5414435223
5355535555 5355555235 5225355535 4545352515 3215122415 2344413314 2313121312 2413414433 2214321
\end{lstlisting}
\end{minipage}\par\smallskip

\noindent\begin{minipage}{\linewidth}\textbf{F1.029}
\begin{lstlisting}
1234122334 4143142213 1213122413 3144432514 2215123515 2535454535 5535225532 5455535555 5355535322
5344145342 2153215132 4122344413 3142213121 3122413414 4332214221 6123412233 4414316221 3151613651
3366463356 6256655666 5643666656 6566666665 6666666666 6666656666 6666666666 6666666666 6666666666
6666666666 6666666666 6666666666 6566666666 6666666566 6666656656 6663465666 5566526653 3646633156
3161513122 6134144332 2143216122 4122334414 3142213121 3122413314 4432214231 5123512243 5414435223
5355535555 5355545235 5225355535 4545352515 3215122415 2344413314 2213121312 2413414433 2214321
\end{lstlisting}
\end{minipage}\par\smallskip

\noindent\begin{minipage}{\linewidth}\textbf{F1.030}
\begin{lstlisting}
1234122334 4143142213 1213122413 3144432514 2315123515 2435454535 5535225532 5555535555 5355535322
5344145342 2153215122 4122344413 3142213121 3122413414 4332214321 6123412233 4414314221 3151613661
3364463356 6236655666 5663666556 6556666666 6666666666 6666656666 6666666666 6666666666 6666666666
6666666666 6666666666 6666666666 6566666666 6666666666 6666556655 6663665666 5566326653 3644633166
3161513122 4134144332 2143216123 4122334414 3142213121 3122413314 4432214221 5123512243 5414435223
5355535555 5355555235 5225355535 4545342515 3215132415 2344413314 2213121312 2413414433 2214321
\end{lstlisting}
\end{minipage}\par\smallskip

\noindent\begin{minipage}{\linewidth}\textbf{F1.031}
\begin{lstlisting}
1234122334 4143142213 1213122413 3144433514 2215123515 2435454535 5535225532 5455335545 3355535322
5344145342 2153215122 4122344413 3142213121 3122413414 4332214321 6123412233 4414314221 3151613641
2366463356 6256655666 5553666656 6566666666 6666666666 6666656666 6666666666 6666666666 6666666666
6666666666 6666666666 6666666666 6566666666 6666666666 6666656656 6663555666 5566526653 3646632146
3161513122 4134144332 2143216123 4122334414 3142213121 3122413314 4432214221 5123512243 5414435223
5355533545 5335545235 5225355535 4545342515 3215122415 3344413314 2213121312 2413414433 2214321
\end{lstlisting}
\end{minipage}\par\smallskip

\noindent\begin{minipage}{\linewidth}\textbf{F1.032}
\begin{lstlisting}
1234122334 4143142213 1213122413 3144433514 2215123515 2535454535 5535225532 5555535545 5355535322
5344145342 2153215132 4122344413 3142213121 3122413414 4332214321 6123412233 4414316221 3151613641
2366463356 6256655666 5653665656 6566666665 6666666666 6666656666 6666666666 6666666666 6666666666
6666666666 6666666666 6666666666 6566666666 6666666566 6666656656 5663565666 5566526653 3646632146
3161513122 6134144332 2143216123 4122334414 3142213121 3122413314 4432214231 5123512243 5414435223
5355535545 5355555235 5225355535 4545352515 3215122415 3344413314 2213121312 2413414433 2214321
\end{lstlisting}
\end{minipage}\par\smallskip

\noindent\begin{minipage}{\linewidth}\textbf{F1.033}
\begin{lstlisting}
1234122334 4143142213 1213122413 3144433514 2215123515 2535454535 5535225533 5555535555 5355535322
5344145342 2153215122 4122344413 3142213121 3122413414 4332214321 6123412233 4414314221 6151613641
3364463356 6256655666 5663666656 6566656666 6666666666 6666656666 6666666666 6666666666 6666666666
6666666666 6666666666 6666666666 6566666666 6666666666 6566656656 6663665666 5566526653 3644633146
3161516122 4134144332 2143216123 4122334414 3142213121 3122413314 4432214221 5123512243 5414435223
5355535555 5355555335 5225355535 4545352515 3215122415 3344413314 2213121312 2413414433 2214321
\end{lstlisting}
\end{minipage}\par\smallskip

\noindent\begin{minipage}{\linewidth}\textbf{F1.034}
\begin{lstlisting}
1234122334 4143142213 1213132413 3144432514 2215123515 2535454535 2535325533 5455535555 5355535322
5344145342 2153215122 4122344413 3142213121 3122413414 4332214321 6123412233 4414314221 3151613651
3324463356 6236655666 5566666656 6566666666 6666666666 6666656666 6666666666 6666666666 6666666666
6666666666 6666666666 6666666666 6566666666 6666666666 6666656656 6666655666 5566326653 3644233156
3161513122 4134144332 2143216123 4122334414 3142213121 3122413314 4432214221 5123512243 5414435223
5355535555 5355545335 5235352535 4545352515 3215122415 2344413314 2313121312 2413414433 2214321
\end{lstlisting}
\end{minipage}\par\smallskip

\noindent\begin{minipage}{\linewidth}\textbf{F1.035}
\begin{lstlisting}
1234122334 4143142213 1213132413 3144432514 2215123515 2535454535 5535225532 5455535555 5355535322
5344145342 2153215122 4122344413 3142213121 3122413414 4332214221 6123412233 4414314221 6151613651
3364463356 6256655666 5663666656 6566656666 6666666666 6666656666 6666666666 6666666666 6666666666
6666666666 6666666666 6666666666 6566666666 6666666666 6566656656 6663665666 5566526653 3644633156
3161516122 4134144332 2143216122 4122334414 3142213121 3122413314 4432214221 5123512243 5414435223
5355535555 5355545235 5225355535 4545352515 3215122415 2344413314 2313121312 2413414433 2214321
\end{lstlisting}
\end{minipage}\par\smallskip

\noindent\begin{minipage}{\linewidth}\textbf{F1.036}
\begin{lstlisting}
1234122334 4143142213 1213132413 3144432514 2315123515 2535454535 5535225532 5455535555 5355535322
5344145342 2153215132 4122344413 3142213121 3122413414 4332214321 6123412233 4414316221 3151613651
3366463356 6256655666 5653666656 6566666665 6666666666 6666656666 6666666666 6666666666 6666666666
6666666666 6666666666 6666666666 6566666666 6666666566 6666656656 6663565666 5566526653 3646633156
3161513122 6134144332 2143216123 4122334414 3142213121 3122413314 4432214231 5123512243 5414435223
5355535555 5355545235 5225355535 4545352515 3215132415 2344413314 2313121312 2413414433 2214321
\end{lstlisting}
\end{minipage}\par\smallskip

\noindent\begin{minipage}{\linewidth}\textbf{F1.037}
\begin{lstlisting}
1234122334 4143142213 1213122413 3144432514 2215123515 2435454535 5535325532 5555335555 5355535322
5344145342 2153215122 4122344413 3142213121 3122413414 4332214321 6123412233 4414314221 3151613641
3366463356 6256655636 5653666653 6566636666 6666663666 6366656666 6666666666 6666666666 6666666666
6666666666 6666666666 6666666666 6566636666 3666666666 6366656356 6663565636 5566526653 3646633146
3161513122 4134144332 2143216123 4122334414 3142213121 3122413314 4432214221 5123512243 5414435223
5355535555 5335555235 5235355535 4545342515 3215122415 2344413314 2213121312 2413414433 2214321
\end{lstlisting}
\end{minipage}\par\smallskip

\noindent\begin{minipage}{\linewidth}\textbf{F1.038}
\begin{lstlisting}
1234122334 4143142213 1213122413 3144432514 2215123515 2535454535 5535325532 5455535555 5355535322
5344145342 2153215122 4122344413 3142213121 3122413414 4332214321 6123412233 4414316221 3151613651
3366463356 6256655666 5653666556 6566656665 6664666666 6666656666 6666666666 6666666666 6666666666
6666666666 6666666666 6666666666 6566666666 6664666566 6566656655 6663565666 5566526653 3646633156
3161513122 6134144332 2143216123 4122334414 3142213121 3122413314 4432214221 5123512243 5414435223
5355535555 5355545235 5235355535 4545352515 3215122415 2344413314 2213121312 2413414433 2214321
\end{lstlisting}
\end{minipage}\par\smallskip

\noindent\begin{minipage}{\linewidth}\textbf{F1.039}
\begin{lstlisting}
1234122334 4143142213 1213122413 3144432514 2215123515 2535454535 5535325533 5555535545 2355235322
5344145342 2153215132 4122344413 3142213121 3122413414 4332214221 6123412233 4414314221 6151613641
3364463356 6256655636 5663665656 6566656666 6666666666 6666656666 6666666666 6666666666 6666666666
6666666666 6666666666 6666666666 6566666666 6666666666 6566656656 5663665636 5566526653 3644633146
3161516122 4134144332 2143216122 4122334414 3142213121 3122413314 4432214231 5123512243 5414435223
5325532545 5355555335 5235355535 4545352515 3215122415 2344413314 2213121312 2413414433 2214321
\end{lstlisting}
\end{minipage}\par\smallskip

\noindent\begin{minipage}{\linewidth}\textbf{F1.040}
\begin{lstlisting}
1234122334 4143142213 1213122413 3144433514 2215123515 2535454535 5535225532 5455335545 3355234322
5344145342 2153215122 4122344413 3142213121 3122413414 4332214321 6123412233 4414314221 3151613641
2266463356 6236655666 5653666656 6566666666 6666663666 6666656666 6666666666 6666666666 6666666666
6666666666 6666666666 6666666666 6566666666 3666666666 6666656656 6663565666 5566326653 3646622146
3161513122 4134144332 2143216123 4122334414 3142213121 3122413314 4432214221 5123512243 5414435223
4325533545 5335545235 5225355535 4545352515 3215122415 3344413314 2213121312 2413414433 2214321
\end{lstlisting}
\end{minipage}\par\smallskip

\noindent\begin{minipage}{\linewidth}\textbf{F1.041}
\begin{lstlisting}
1234122334 4143142213 1213122413 3144433514 2215123515 2535454535 5535225533 5555325555 5355235322
5344145342 2153215132 4122344413 3142213121 3122413414 4332214321 6123412233 4414314221 3151613641
3266463356 6236655636 5643666656 6566666666 6666666666 6666656666 6666666666 6666666666 6666666666
6666666666 6666666666 6666666666 6566666666 6666666666 6666656656 6663465636 5566326653 3646623146
3161513122 4134144332 2143216123 4122334414 3142213121 3122413314 4432214231 5123512243 5414435223
5325535555 5235555335 5225355535 4545352515 3215122415 3344413314 2213121312 2413414433 2214321
\end{lstlisting}
\end{minipage}\par\smallskip

\noindent\begin{minipage}{\linewidth}\textbf{F1.042}
\begin{lstlisting}
1234122334 4143142213 1213132413 3144432514 2215123515 2435454535 5535325533 5555335555 5355535322
5344145342 2153215132 4122344413 3142213121 3122413414 4332214321 6123412233 4414314221 3151613641
3326463356 6256655666 5656665656 6566666666 6666646666 6666656666 6666666666 6666666666 6666666666
6666666666 6666666666 6666666666 6566666666 6466666666 6666656656 5666565666 5566526653 3646233146
3161513122 4134144332 2143216123 4122334414 3142213121 3122413314 4432214231 5123512243 5414435223
5355535555 5335555335 5235355535 4545342515 3215122415 2344413314 2313121312 2413414433 2214321
\end{lstlisting}
\end{minipage}\par\smallskip

\noindent\begin{minipage}{\linewidth}\textbf{F1.043}
\begin{lstlisting}
1234122334 4143142213 1213132413 3144432514 2215123515 2535454535 5535225532 5555535555 5355534322
5344145342 2153215122 4122344413 3142213121 3122413414 4332214221 6123412233 4414314221 6151613661
3364463356 6256655666 5663666556 6566656666 6666666666 6666656666 6666666666 6666666666 6666666666
6666666666 6666666666 6666666666 6566666666 6666666666 6566656655 6663665666 5566526653 3644633166
3161516122 4134144332 2143216122 4122334414 3142213121 3122413314 4432214221 5123512243 5414435223
4355535555 5355555235 5225355535 4545352515 3215122415 2344413314 2313121312 2413414433 2214321
\end{lstlisting}
\end{minipage}\par\smallskip

\noindent\begin{minipage}{\linewidth}\textbf{F1.044}
\begin{lstlisting}
1234122334 4143142213 1213132413 3144432514 2215123515 2535454535 5535325533 5555535545 2355535322
5344145342 2153215122 4122344413 3142213121 3122413414 4332214321 6123412233 4416316221 3151613641
3324463356 6256655636 5663665653 6566666665 6646666666 6666656666 6666666666 6666666666 6666666666
6666666666 6666666666 6666666666 6566666666 6666466566 6666656356 5663665636 5566526653 3644233146
3161513122 6136144332 2143216123 4122334414 3142213121 3122413314 4432214221 5123512243 5414435223
5355532545 5355555335 5235355535 4545352515 3215122415 2344413314 2313121312 2413414433 2214321
\end{lstlisting}
\end{minipage}\par\smallskip

\noindent\begin{minipage}{\linewidth}\textbf{F1.045}
\begin{lstlisting}
1234122334 4143142213 1213132413 3144432514 2315123515 2435454535 2535225532 5455335555 2355235322
5344145342 2153215122 4122344413 3142213121 3122413414 4332214321 6123412233 4414314221 3151613641
3364463356 6236655636 5643666453 6566666666 6666666666 6666656666 6666666666 6666666666 6666666666
6666666666 6666666666 6666666666 6566666666 6666666666 6666656354 6663465636 5566326653 3644633146
3161513122 4134144332 2143216123 4122334414 3142213121 3122413314 4432214221 5123512243 5414435223
5325532555 5335545235 5225352535 4545342515 3215132415 2344413314 2313121312 2413414433 2214321
\end{lstlisting}
\end{minipage}\par\smallskip

\noindent\begin{minipage}{\linewidth}\textbf{F1.046}
\begin{lstlisting}
1234122334 4143142213 1213132413 3144432514 2315123515 2535454535 2535225533 5555535555 5355535322
5344145342 2153215132 4122344413 3142213121 3122413414 4332214221 6123412233 4414314221 6151613641
3326463356 6236655666 5646666656 6566636666 6666666666 6666656666 6666666666 6666666666 6666666666
6666666666 6666666666 6666666666 6566666666 6666666666 6366656656 6666465666 5566326653 3646233146
3161516122 4134144332 2143216122 4122334414 3142213121 3122413314 4432214231 5123512243 5414435223
5355535555 5355555335 5225352535 4545352515 3215132415 2344413314 2313121312 2413414433 2214321
\end{lstlisting}
\end{minipage}\par\smallskip

\noindent\begin{minipage}{\linewidth}\textbf{F1.047}
\begin{lstlisting}
1234122334 4143142213 1213122413 3144432514 2215123515 2435454535 5535225533 5555335555 5355535322
5344145342 2153215132 4122344413 3142213121 3122413414 4332214321 6123412233 4414316221 3151613641
3366463356 6256655666 5643665656 6556666665 6666666666 6666656666 6666666666 6666666666 6666666666
6666666666 6666666666 6666666666 6566666666 6666666566 6666556656 5663465666 5566526653 3646633146
3161513122 6134144332 2143216123 4122334414 3142213121 3122413314 4432214231 5123512243 5414435223
5355535555 5335555335 5225355535 4545342515 3215122415 2344413314 2213121312 2413414433 2214321
\end{lstlisting}
\end{minipage}\par\smallskip

\noindent\begin{minipage}{\linewidth}\textbf{F1.048}
\begin{lstlisting}
1234122334 4143142213 1213122413 3144432514 2215123515 2435454535 5535325533 5455535555 2355535322
5344145342 2153215132 4122344413 3142213121 3122413414 4332214221 6123412233 4414314221 6151613661
3364463356 6236655666 5663666556 6566656666 6666666666 6666656666 6666666666 6666666666 6666666666
6666666666 6666666666 6666666666 6566666666 6666666666 6566656655 6663665666 5566326653 3644633166
3161516122 4134144332 2143216122 4122334414 3142213121 3122413314 4432214231 5123512243 5414435223
5355532555 5355545335 5235355535 4545342515 3215122415 2344413314 2213121312 2413414433 2214321
\end{lstlisting}
\end{minipage}\par\smallskip

\noindent\begin{minipage}{\linewidth}\textbf{F1.049}
\begin{lstlisting}
1234122334 4143142213 1213122413 3144432514 2215123515 2535454535 5535225533 5555335555 5355235322
5344145342 2153215122 4122344413 3142213121 3122413414 4332214321 6123412233 4414316221 3151613661
3364463356 6236655666 5663666556 6566666665 6666666666 6666656666 6666666666 6666666666 6666666666
6666666666 6666666666 6666666666 6566666666 6666666566 6666656655 6663665666 5566326653 3644633166
3161513122 6134144332 2143216123 4122334414 3142213121 3122413314 4432214221 5123512243 5414435223
5325535555 5335555335 5225355535 4545352515 3215122415 2344413314 2213121312 2413414433 2214321
\end{lstlisting}
\end{minipage}\par\smallskip

\noindent\begin{minipage}{\linewidth}\textbf{F1.050}
\begin{lstlisting}
1234122334 4143142213 1213122413 3144432514 2315123515 2535454535 2535325532 5455535545 5355235322
5344145342 2153215132 4122344413 3142213121 3122413414 4332214321 6123412233 4414314221 6151613661
3364463356 6236655666 5663666556 6566656666 6666666666 6666656666 6666666666 6666666666 6666666666
6666666666 6666666666 6666666666 6566666666 6666666666 6566656655 6663665666 5566326653 3644633166
3161516122 4134144332 2143216123 4122334414 3142213121 3122413314 4432214231 5123512243 5414435223
5325535545 5355545235 5235352535 4545352515 3215132415 2344413314 2213121312 2413414433 2214321
\end{lstlisting}
\end{minipage}\par\smallskip

\noindent\begin{minipage}{\linewidth}\textbf{F1.051}
\begin{lstlisting}
1234122334 4143142213 1213122413 3144432514 2315123515 2535454535 5535225532 5455335555 5355235322
5344145342 2153215122 4122344413 3142213121 3122413414 4332214321 6123412233 4414314221 6151613641
3364463356 6256655666 5663666656 6566636666 6666666666 6666656666 6666666666 6666666666 6666666666
6666666666 6666666666 6666666666 6566666666 6666666666 6366656656 6663665666 5566526653 3644633146
3161516122 4134144332 2143216123 4122334414 3142213121 3122413314 4432214221 5123512243 5414435223
5325535555 5335545235 5225355535 4545352515 3215132415 2344413314 2213121312 2413414433 2214321
\end{lstlisting}
\end{minipage}\par\smallskip

\noindent\begin{minipage}{\linewidth}\textbf{F1.052}
\begin{lstlisting}
1234122334 4143142213 1213122413 3144433514 2215123515 2535454535 2535325532 5555535555 5355535322
5344145342 2153215122 4122344413 3142213121 3122413414 4332214321 6123412233 4414314221 6151613661
3264463356 6256655666 5663666456 6566656666 6666666666 6666656666 6666666666 6666666666 6666666666
6666666666 6666666666 6666666666 6566666666 6666666666 6566656654 6663665666 5566526653 3644623166
3161516122 4134144332 2143216123 4122334414 3142213121 3122413314 4432214221 5123512243 5414435223
5355535555 5355555235 5235352535 4545352515 3215122415 3344413314 2213121312 2413414433 2214321
\end{lstlisting}
\end{minipage}\par\smallskip

\noindent\begin{minipage}{\linewidth}\textbf{F1.053}
\begin{lstlisting}
1234122334 4143142213 1213122413 3144433514 2215123515 2535454535 5535225532 5555335555 5355535322
5344145342 2153215122 4122344413 3142213121 3122413414 4332214321 6123412233 4414314221 6151613661
3264463356 6236655666 5663666556 6566656666 6666666666 6666656666 6666666666 6666666666 6666666666
6666666666 6666666666 6666666666 6566666666 6666666666 6566656655 6663665666 5566326653 3644623166
3161516122 4134144332 2143216123 4122334414 3142213121 3122413314 4432214221 5123512243 5414435223
5355535555 5335555235 5225355535 4545352515 3215122415 3344413314 2213121312 2413414433 2214321
\end{lstlisting}
\end{minipage}\par\smallskip

\noindent\begin{minipage}{\linewidth}\textbf{F1.054}
\begin{lstlisting}
1234122334 4143142213 1213122413 3144433514 2215123515 2535454535 5535325532 5555535555 5355535322
5344145342 2153215122 4122344413 3142213121 3122413414 4332214321 6123412233 4414314221 6151613661
2264463356 6256655666 5663666556 6566656666 6666666666 6666656666 6666666666 6666666666 6666666666
6666666666 6666666666 6666666666 6566666666 6666666666 6566656655 6663665666 5566526653 3644622166
3161516122 4134144332 2143216123 4122334414 3142213121 3122413314 4432214221 5123512243 5414435223
5355535555 5355555235 5235355535 4545352515 3215122415 3344413314 2213121312 2413414433 2214321
\end{lstlisting}
\end{minipage}\par\smallskip

\noindent\begin{minipage}{\linewidth}\textbf{F1.055}
\begin{lstlisting}
1234122334 4143142213 1213132413 3144432514 2215123515 2435454535 2535325533 5555235545 5355235322
5344145342 2153215132 4122344413 3142213121 3122413414 4332214321 6123412233 4414314221 6151613661
3364463356 6236655666 5663666556 6566656666 6666666666 6666656666 6666666666 6666666666 6666666666
6666666666 6666666666 6666666666 6566666666 6666666666 6566656655 6663665666 5566326653 3644633166
3161516122 4134144332 2143216123 4122334414 3142213121 3122413314 4432214231 5123512243 5414435223
5325535545 5325555335 5235352535 4545342515 3215122415 2344413314 2313121312 2413414433 2214321
\end{lstlisting}
\end{minipage}\par\smallskip

\noindent\begin{minipage}{\linewidth}\textbf{F1.056}
\begin{lstlisting}
1234122334 4143142213 1213132413 3144432514 2215123515 2435454535 5535325533 5555535555 5355535322
5344145342 2153215132 4122344413 3142213121 3122413414 4332214221 6123412233 4414316221 3151613661
3364463356 6256655666 5663666456 6566666665 6666666666 6666656666 6666666666 6666666666 6666666666
6666666666 6666666666 6666666666 6566666666 6666666566 6666656654 6663665666 5566526653 3644633166
3161513122 6134144332 2143216122 4122334414 3142213121 3122413314 4432214231 5123512243 5414435223
5355535555 5355555335 5235355535 4545342515 3215122415 2344413314 2313121312 2413414433 2214321
\end{lstlisting}
\end{minipage}\par\smallskip

\noindent\begin{minipage}{\linewidth}\textbf{F1.057}
\begin{lstlisting}
1234122334 4143142213 1213132413 3144432514 2315123515 2535454535 5535225532 5555335555 5355535322
5344145342 2153215122 4122344413 3142213121 3122413414 4332214221 6123412233 4414314221 6151613661
3364463356 6256655666 5663666556 6566656666 6666666666 6666656666 6666666666 6666666666 6666666666
6666666666 6666666666 6666666666 6566666666 6666666666 6566656655 6663665666 5566526653 3644633166
3161516122 4134144332 2143216122 4122334414 3142213121 3122413314 4432214221 5123512243 5414435223
5355535555 5335555235 5225355535 4545352515 3215132415 2344413314 2313121312 2413414433 2214321
\end{lstlisting}
\end{minipage}\par\smallskip

\noindent\begin{minipage}{\linewidth}\textbf{F1.058}
\begin{lstlisting}
1234122334 4143142213 1213132413 3144432514 2315123515 2535454535 5535325533 5555535545 2355535322
5344145342 2153215122 4122344413 3142213121 3122413414 4332214321 6123412233 4416316221 3151613641
3324463356 6256655636 5663665653 6566666665 6646666666 6666656666 6666666666 6666666666 6666666666
6666666666 6666666666 6666666666 6566666666 6666466566 6666656356 5663665636 5566526653 3644233146
3161513122 6136144332 2143216123 4122334414 3142213121 3122413314 4432214221 5123512243 5414435223
5355532545 5355555335 5235355535 4545352515 3215132415 2344413314 2313121312 2413414433 2214321
\end{lstlisting}
\end{minipage}\par\smallskip

\noindent\begin{minipage}{\linewidth}\textbf{F1.059}
\begin{lstlisting}
1234122334 4143142213 1213132413 3144432514 2315123515 2535454535 5535325533 5555535555 3355535322
5344145342 2153215122 4122344413 3142213121 3122413414 4332214321 6123412233 4416314221 6151613641
3324463356 6236655636 5646666653 6566636666 6646666666 6666656666 6666666666 6666666666 6666666666
6666666666 6666666666 6666666666 6566666666 6666466666 6366656356 6666465636 5566326653 3644233146
3161516122 4136144332 2143216123 4122334414 3142213121 3122413314 4432214221 5123512243 5414435223
5355533555 5355555335 5235355535 4545352515 3215132415 2344413314 2313121312 2413414433 2214321
\end{lstlisting}
\end{minipage}\par\smallskip

\noindent\begin{minipage}{\linewidth}\textbf{F1.060}
\begin{lstlisting}
1234122334 4143142213 1213122413 3144432514 2215123515 2435454535 2535225532 5555535555 5355235322
5344145342 2153215132 4122344413 3142213121 3122413414 4332214321 6123412233 4414316221 3151613641
3364463356 6256655636 5663666656 6566666665 6666666666 6666656666 6666666666 6666666666 6666666666
6666666666 6666666666 6666666666 6566666666 6666666566 6666656656 6663665636 5566526653 3644633146
3161513122 6134144332 2143216123 4122334414 3142213121 3122413314 4432214231 5123512243 5414435223
5325535555 5355555235 5225352535 4545342515 3215122415 2344413314 2213121312 2413414433 2214321
\end{lstlisting}
\end{minipage}\par\smallskip

\noindent\begin{minipage}{\linewidth}\textbf{F1.061}
\begin{lstlisting}
1234122334 4143142213 1213122413 3144432514 2215123515 2435454535 2535225532 5555535555 5355535322
5344145342 2153215122 4122344413 3142213121 3122413414 4332214221 6123412233 4414316221 3151613661
3364463356 6236655666 5663666453 6566666665 6666666666 6666656666 6666666666 6666666666 6666666666
6666666666 6666666666 6666666666 6566666666 6666666566 6666656354 6663665666 5566326653 3644633166
3161513122 6134144332 2143216122 4122334414 3142213121 3122413314 4432214221 5123512243 5414435223
5355535555 5355555235 5225352535 4545342515 3215122415 2344413314 2213121312 2413414433 2214321
\end{lstlisting}
\end{minipage}\par\smallskip

\noindent\begin{minipage}{\linewidth}\textbf{F1.062}
\begin{lstlisting}
1234122334 4143142213 1213122413 3144432514 2215123515 2535454535 5535325533 5555535555 5355234322
5344145342 2153215132 4122344413 3142213121 3122413414 4332214321 6123412233 4414314221 3151613641
3366463356 6236655666 5653665556 6566566666 3646663666 6666656666 6666666666 6666666666 6666666666
6666666666 6666666666 6666666666 6566666666 3666463666 6656656655 5663565666 5566326653 3646633146
3161513122 4134144332 2143216123 4122334414 3142213121 3122413314 4432214231 5123512243 5414435223
4325535555 5355555335 5235355535 4545352515 3215122415 2344413314 2213121312 2413414433 2214321
\end{lstlisting}
\end{minipage}\par\smallskip

\noindent\begin{minipage}{\linewidth}\textbf{F1.063}
\begin{lstlisting}
1234122334 4143142213 1213122413 3144432514 2315123515 2535454535 2535225533 5455335545 3355234322
5344145342 2153215122 4122344413 3142213121 3122413414 4332214321 6123412233 4414314221 3151613641
3364463356 6236655636 5663666556 6566666666 6666666666 6666656666 6666666666 6666666666 6666666666
6666666666 6666666666 6666666666 6566666666 6666666666 6666656655 6663665636 5566326653 3644633146
3161513122 4134144332 2143216123 4122334414 3142213121 3122413314 4432214221 5123512243 5414435223
4325533545 5335545335 5225352535 4545352515 3215132415 2344413314 2213121312 2413414433 2214321
\end{lstlisting}
\end{minipage}\par\smallskip

\noindent\begin{minipage}{\linewidth}\textbf{F1.064}
\begin{lstlisting}
1234122334 4143142213 1213122413 3144432514 2315123515 2535454535 5535325533 5555535555 5355535322
5344145342 2153215132 4122344413 3142213121 3122413414 4332214221 6123412233 4414314221 6151613661
3364463356 6256655666 5663666556 6566656666 6666666666 6666656666 6666666666 6666666666 6666666666
6666666666 6666666666 6666666666 6566666666 6666666666 6566656655 6663665666 5566526653 3644633166
3161516122 4134144332 2143216122 4122334414 3142213121 3122413314 4432214231 5123512243 5414435223
5355535555 5355555335 5235355535 4545352515 3215132415 2344413314 2213121312 2413414433 2214321
\end{lstlisting}
\end{minipage}\par\smallskip

\noindent\begin{minipage}{\linewidth}\textbf{F1.065}
\begin{lstlisting}
1234122334 4143142213 1213122413 3144433514 2215123515 2435454535 2535325533 5555335545 2355235322
5344145342 2153215122 4122344413 3142213121 3122413414 4332214321 6123412233 4414314221 3151613641
3364463356 6236655666 5653666556 6566656665 6666663666 6636656666 6666666666 6666666666 6666666666
6666666666 6666666666 6666666666 6566366666 3666666566 6566656655 6663565666 5566326653 3644633146
3161513122 4134144332 2143216123 4122334414 3142213121 3122413314 4432214221 5123512243 5414435223
5325532545 5335555335 5235352535 4545342515 3215122415 3344413314 2213121312 2413414433 2214321
\end{lstlisting}
\end{minipage}\par\smallskip

\noindent\begin{minipage}{\linewidth}\textbf{F1.066}
\begin{lstlisting}
1234122334 4143142213 1213122413 3144433514 2215123515 2435454535 5535225533 5555325545 5355535322
5344145342 2153215132 4122344413 3142213121 3122413414 4332214321 6123412233 4414314221 3151613661
3264463356 6236655666 5663666456 6566666666 6666666666 6666656666 6666666666 6666666666 6666666666
6666666666 6666666666 6666666666 6566666666 6666666666 6666656654 6663665666 5566326653 3644623166
3161513122 4134144332 2143216123 4122334414 3142213121 3122413314 4432214231 5123512243 5414435223
5355535545 5235555335 5225355535 4545342515 3215122415 3344413314 2213121312 2413414433 2214321
\end{lstlisting}
\end{minipage}\par\smallskip

\noindent\begin{minipage}{\linewidth}\textbf{F1.067}
\begin{lstlisting}
1234122334 4143142213 1213122413 3144433514 2215123515 2435454535 5535325532 5555535555 5355235322
5344145342 2153215122 4122344413 3142213121 3122413414 4332214321 6123412233 4414314221 6151613661
3264463356 6256655666 5663666456 6566656666 6666666666 6666656666 6666666666 6666666666 6666666666
6666666666 6666666666 6666666666 6566666666 6666666666 6566656654 6663665666 5566526653 3644623166
3161516122 4134144332 2143216123 4122334414 3142213121 3122413314 4432214221 5123512243 5414435223
5325535555 5355555235 5235355535 4545342515 3215122415 3344413314 2213121312 2413414433 2214321
\end{lstlisting}
\end{minipage}\par\smallskip

\noindent\begin{minipage}{\linewidth}\textbf{F1.068}
\begin{lstlisting}
1234122334 4143142213 1213122413 3144433514 2215123515 2435454535 5535325533 5455535555 5355535322
5344145342 2153215122 4122344413 3142213121 3122413414 4332214321 6123412233 4414316221 3151613651
2266463356 6256655666 5653666656 6566666665 6666666666 6666656666 6666666666 6666666666 6666666666
6666666666 6666666666 6666666666 6566666666 6666666566 6666656656 6663565666 5566526653 3646622156
3161513122 6134144332 2143216123 4122334414 3142213121 3122413314 4432214221 5123512243 5414435223
5355535555 5355545335 5235355535 4545342515 3215122415 3344413314 2213121312 2413414433 2214321
\end{lstlisting}
\end{minipage}\par\smallskip

\noindent\begin{minipage}{\linewidth}\textbf{F1.069}
\begin{lstlisting}
1234122334 4143142213 1213122413 3144433514 2315123515 2535454535 5535325533 5455525545 5355235322
5344145342 2153215122 4122344413 3142213121 3122413414 4332214321 6123412233 4414314221 6151613661
3364463356 6236655666 5663666556 6566656666 6666666666 6666656666 6666666666 6666666666 6666666666
6666666666 6666666666 6666666666 6566666666 6666666666 6566656655 6663665666 5566326653 3644633166
3161516122 4134144332 2143216123 4122334414 3142213121 3122413314 4432214221 5123512243 5414435223
5325535545 5255545335 5235355535 4545352515 3215132415 3344413314 2213121312 2413414433 2214321
\end{lstlisting}
\end{minipage}\par\smallskip

\noindent\begin{minipage}{\linewidth}\textbf{F1.070}
\begin{lstlisting}
1234122334 4143142213 1213132413 3144432514 2215123515 2535454535 5535225532 5455335555 5355535322
5344145342 2153215132 4122344413 3142213121 3122413414 4332214321 6123412233 4414314221 6151613661
3364463356 6256655666 5663665556 6566656666 6666666666 6666656666 6666666666 6666666666 6666666666
6666666666 6666666666 6666666666 6566666666 6666666666 6566656655 5663665666 5566526653 3644633166
3161516122 4134144332 2143216123 4122334414 3142213121 3122413314 4432214231 5123512243 5414435223
5355535555 5335545235 5225355535 4545352515 3215122415 2344413314 2313121312 2413414433 2214321
\end{lstlisting}
\end{minipage}\par\smallskip

\noindent\begin{minipage}{\linewidth}\textbf{F1.071}
\begin{lstlisting}
1234122334 4143142213 1213132413 3144432514 2215123515 2535454535 5535325532 5555535555 5355534322
5344145342 2153215122 4122344413 3142213121 3122413414 4332214221 6123412233 4414316221 3151613641
3366463356 6256655666 5653666453 6566636665 6646663666 6636656666 6666666666 6666666666 6666666666
6666666666 6666666666 6666666666 6566366666 3666466566 6366656354 6663565666 5566526653 3646633146
3161513122 6134144332 2143216122 4122334414 3142213121 3122413314 4432214221 5123512243 5414435223
4355535555 5355555235 5235355535 4545352515 3215122415 2344413314 2313121312 2413414433 2214321
\end{lstlisting}
\end{minipage}\par\smallskip

\noindent\begin{minipage}{\linewidth}\textbf{F1.072}
\begin{lstlisting}
1234122334 4143142213 1213132413 3144432514 2215123515 2535454535 5535325533 5455535545 2355235322
5344145342 2153215132 4122344413 3142213121 3122413414 4332214221 6123412233 4416314221 6151613641
3324463356 6236655636 5646666653 6566636666 6646666666 6666656666 6666666666 6666666666 6666666666
6666666666 6666666666 6666666666 6566666666 6666466666 6366656356 6666465636 5566326653 3644233146
3161516122 4136144332 2143216122 4122334414 3142213121 3122413314 4432214231 5123512243 5414435223
5325532545 5355545335 5235355535 4545352515 3215122415 2344413314 2313121312 2413414433 2214321
\end{lstlisting}
\end{minipage}\par\smallskip

\noindent\begin{minipage}{\linewidth}\textbf{F1.073}
\begin{lstlisting}
1234122334 4143142213 1213132413 3144432514 2315123515 2535454535 2535225532 5455535555 3355534322
5344145342 2153215132 4122344413 3142213121 3122413414 4332214321 6123412233 4414316221 3151613651
3366463356 6256655666 5653666656 6566666665 6664666666 6666656666 6666666666 6666666666 6666666666
6666666666 6666666666 6666666666 6566666666 6664666566 6666656656 6663565666 5566526653 3646633156
3161513122 6134144332 2143216123 4122334414 3142213121 3122413314 4432214231 5123512243 5414435223
4355533555 5355545235 5225352535 4545352515 3215132415 2344413314 2313121312 2413414433 2214321
\end{lstlisting}
\end{minipage}\par\smallskip

\noindent\begin{minipage}{\linewidth}\textbf{F1.074}
\begin{lstlisting}
1234122334 4143142213 1213132413 3144432514 2315123515 2535454535 5535325533 5555535555 2355535322
5344145342 2153215122 4122344413 3142213121 3122413414 4332214321 6123412233 4416316221 3151613641
3324463356 6256655636 5663665653 6566666665 6646666666 6666656666 6666666666 6666666666 6666666666
6666666666 6666666666 6666666666 6566666666 6666466566 6666656356 5663665636 5566526653 3644233146
3161513122 6136144332 2143216123 4122334414 3142213121 3122413314 4432214221 5123512243 5414435223
5355532555 5355555335 5235355535 4545352515 3215132415 2344413314 2313121312 2413414433 2214321
\end{lstlisting}
\end{minipage}\par\smallskip

\noindent\begin{minipage}{\linewidth}\textbf{F1.075}
\begin{lstlisting}
1234122334 4143142213 1213122413 3144432514 2215123515 2535454535 5535225533 5555335555 5355534322
5344145342 2153215132 4122344413 3142213121 3122413414 4332214321 6123412233 4414314221 6151613661
3364463356 6256655666 5663666556 6566656666 6666666666 6666656666 6666666666 6666666666 6666666666
6666666666 6666666666 6666666666 6566666666 6666666666 6566656655 6663665666 5566526653 3644633166
3161516122 4134144332 2143216123 4122334414 3142213121 3122413314 4432214231 5123512243 5414435223
4355535555 5335555335 5225355535 4545352515 3215122415 2344413314 2213121312 2413414433 2214321
\end{lstlisting}
\end{minipage}\par\smallskip

\noindent\begin{minipage}{\linewidth}\textbf{F1.076}
\begin{lstlisting}
1234122334 4143142213 1213122413 3144432514 2315123515 2435454535 2535225533 5555335545 3355535322
5344145342 2153215122 4122344413 3142213121 3122413414 4332214321 6123412233 4414314221 3151613661
3364463356 6256655666 5663666456 6566666666 3666666666 6636656666 6666666666 6666666666 6666666666
6666666666 6666666666 6666666666 6566366666 6666663666 6666656654 6663665666 5566526653 3644633166
3161513122 4134144332 2143216123 4122334414 3142213121 3122413314 4432214221 5123512243 5414435223
5355533545 5335555335 5225352535 4545342515 3215132415 2344413314 2213121312 2413414433 2214321
\end{lstlisting}
\end{minipage}\par\smallskip

\noindent\begin{minipage}{\linewidth}\textbf{F1.077}
\begin{lstlisting}
1234122334 4143142213 1213122413 3144433514 2215123515 2535454535 5535225533 5555535545 2355535322
5344145342 2153215132 4122344413 3142213121 3122413414 4332214321 6123412233 4414316221 3151613661
3364463356 6236655636 5553666456 6556666665 6646666666 6666656666 6666666666 6666666666 6666666666
6666666666 6666666666 6666666666 6566666666 6666466566 6666556654 6663555636 5566326653 3644633166
3161513122 6134144332 2143216123 4122334414 3142213121 3122413314 4432214231 5123512243 5414435223
5355532545 5355555335 5225355535 4545352515 3215122415 3344413314 2213121312 2413414433 2214321
\end{lstlisting}
\end{minipage}\par\smallskip

\noindent\begin{minipage}{\linewidth}\textbf{F1.078}
\begin{lstlisting}
1234122334 4143142213 1213122413 3144433514 2315123515 2535454535 5535225533 5455535555 5355534322
5344145342 2153215122 4122344413 3142213121 3122413414 4332214321 6123412233 4414316221 3151613651
2366463356 6256655666 5653666656 6566666665 6666666666 6666656666 6666666666 6666666666 6666666666
6666666666 6666666666 6666666666 6566666666 6666666566 6666656656 6663565666 5566526653 3646632156
3161513122 6134144332 2143216123 4122334414 3142213121 3122413314 4432214221 5123512243 5414435223
4355535555 5355545335 5225355535 4545352515 3215132415 3344413314 2213121312 2413414433 2214321
\end{lstlisting}
\end{minipage}\par\smallskip

\noindent\begin{minipage}{\linewidth}\textbf{F1.079}
\begin{lstlisting}
1234122334 4143142213 1213132413 3144432514 2315123515 2535454535 5535325533 5555535555 2355535322
5344145342 2153215132 4122344413 3142213121 3122413414 4332214321 6123412233 4416316221 3151613641
3324463356 6256655636 5663665653 6566666665 6646666666 6666656666 6666666666 6666666666 6666666666
6666666666 6666666666 6666666666 6566666666 6666466566 6666656356 5663665636 5566526653 3644233146
3161513122 6136144332 2143216123 4122334414 3142213121 3122413314 4432214231 5123512243 5414435223
5355532555 5355555335 5235355535 4545352515 3215132415 2344413314 2313121312 2413414433 2214321
\end{lstlisting}
\end{minipage}\par\smallskip

\noindent\begin{minipage}{\linewidth}\textbf{F1.080}
\begin{lstlisting}
1234122334 4143142213 1213122413 3144432514 2215123515 2435454535 2535325533 5455535545 3355235322
5344145342 2153215122 4122344413 3142213121 3122413414 4332214321 6123412233 4414316221 3151613651
3366463356 6236655666 5653666653 6556636665 6664666666 6636656666 6666666666 6666666666 6666666666
6666666666 6666666666 6666666666 6566366666 6664666566 6366556356 6663565666 5566326653 3646633156
3161513122 6134144332 2143216123 4122334414 3142213121 3122413314 4432214221 5123512243 5414435223
5325533545 5355545335 5235352535 4545342515 3215122415 2344413314 2213121312 2413414433 2214321
\end{lstlisting}
\end{minipage}\par\smallskip

\noindent\begin{minipage}{\linewidth}\textbf{F1.081}
\begin{lstlisting}
1234122334 4143142213 1213122413 3144432514 2215123515 2435454535 5535225532 5555335555 5355535322
5344145342 2153215122 4122344413 3142213121 3122413414 4332214321 6123412233 4414314221 6151613641
3364463356 6256655636 5543665456 6566636665 6666646666 6666656666 6666666666 6666666666 6666666666
6666666666 6666666666 6666666666 6566666666 6466666566 6366656654 5663455636 5566526653 3644633146
3161516122 4134144332 2143216123 4122334414 3142213121 3122413314 4432214221 5123512243 5414435223
5355535555 5335555235 5225355535 4545342515 3215122415 2344413314 2213121312 2413414433 2214321
\end{lstlisting}
\end{minipage}\par\smallskip

\noindent\begin{minipage}{\linewidth}\textbf{F1.082}
\begin{lstlisting}
1234122334 4143142213 1213122413 3144432514 2215123515 2435454535 5535225533 5455335555 5355535322
5344145342 2153215132 4122344413 3142213121 3122413414 4332214221 6123412233 4414314221 6151613661
3364463356 6236655666 5663666456 6566636666 6666666666 6666656666 6666666666 6666666666 6666666666
6666666666 6666666666 6666666666 6566666666 6666666666 6366656654 6663665666 5566326653 3644633166
3161516122 4134144332 2143216122 4122334414 3142213121 3122413314 4432214231 5123512243 5414435223
5355535555 5335545335 5225355535 4545342515 3215122415 2344413314 2213121312 2413414433 2214321
\end{lstlisting}
\end{minipage}\par\smallskip

\noindent\begin{minipage}{\linewidth}\textbf{F1.083}
\begin{lstlisting}
1234122334 4143142213 1213122413 3144432514 2215123515 2535454535 5535225532 5555335555 2355235322
5344145342 2153215132 4122344413 3142213121 3122413414 4332214321 6123412233 4414316221 3151613641
3366463356 6236655636 5553666656 6566636665 3666643666 6636656666 6666666666 6666666666 6666666666
6666666666 6666666666 6666666666 6566366666 3466663566 6366656656 6663555636 5566326653 3646633146
3161513122 6134144332 2143216123 4122334414 3142213121 3122413314 4432214231 5123512243 5414435223
5325532555 5335555235 5225355535 4545352515 3215122415 2344413314 2213121312 2413414433 2214321
\end{lstlisting}
\end{minipage}\par\smallskip

\noindent\begin{minipage}{\linewidth}\textbf{F1.084}
\begin{lstlisting}
1234122334 4143142213 1213122413 3144433514 2215123515 2435454535 5535225532 5555335555 3355535322
5344145342 2153215122 4122344413 3142213121 3122413414 4332214321 6123412233 4414316221 3151613661
2264463356 6256655666 5663666556 6566666665 6666666666 6666656666 6666666666 6666666666 6666666666
6666666666 6666666666 6666666666 6566666666 6666666566 6666656655 6663665666 5566526653 3644622166
3161513122 6134144332 2143216123 4122334414 3142213121 3122413314 4432214221 5123512243 5414435223
5355533555 5335555235 5225355535 4545342515 3215122415 3344413314 2213121312 2413414433 2214321
\end{lstlisting}
\end{minipage}\par\smallskip

\noindent\begin{minipage}{\linewidth}\textbf{F1.085}
\begin{lstlisting}
1234122334 4143142213 1213122413 3144433514 2215123515 2435454535 5535225533 5555535555 3355535322
5344145342 2153215132 4122344413 3142213121 3122413414 4332214321 6123412233 4414316221 3151613641
2266463356 6256655636 5653666656 6566666665 6666666666 6636656666 6666666666 6666666666 6666666666
6666666666 6666666666 6666666666 6566366666 6666666566 6666656656 6663565636 5566526653 3646622146
3161513122 6134144332 2143216123 4122334414 3142213121 3122413314 4432214231 5123512243 5414435223
5355533555 5355555335 5225355535 4545342515 3215122415 3344413314 2213121312 2413414433 2214321
\end{lstlisting}
\end{minipage}\par\smallskip

\noindent\begin{minipage}{\linewidth}\textbf{F1.086}
\begin{lstlisting}
1234122334 4143142213 1213122413 3144433514 2315123515 2535454535 5535325533 5555325555 5355234322
5344145342 2153215122 4122344413 3142213121 3122413414 4332214321 6123412233 4414316221 3151613641
3366463356 6256655666 5553666653 6566666665 6666666666 6666656666 6666666666 6666666666 6666666666
6666666666 6666666666 6666666666 6566666666 6666666566 6666656356 6663555666 5566526653 3646633146
3161513122 6134144332 2143216123 4122334414 3142213121 3122413314 4432214221 5123512243 5414435223
4325535555 5235555335 5235355535 4545352515 3215132415 3344413314 2213121312 2413414433 2214321
\end{lstlisting}
\end{minipage}\par\smallskip

\noindent\begin{minipage}{\linewidth}\textbf{F1.087}
\begin{lstlisting}
1234122334 4143142213 1213132413 3144432514 2215123515 2435454535 5535325532 5455335545 2355235322
5344145342 2153215122 4122344413 3142213121 3122413414 4332214321 6123412233 4414314221 3151613641
3364463356 6236655636 5643665453 6556666665 3646663666 6336656666 6666666666 6666666666 6666666666
6666666666 6666666666 6666666666 6566336666 3666463566 6666556354 5663465636 5566326653 3644633146
3161513122 4134144332 2143216123 4122334414 3142213121 3122413314 4432214221 5123512243 5414435223
5325532545 5335545235 5235355535 4545342515 3215122415 2344413314 2313121312 2413414433 2214321
\end{lstlisting}
\end{minipage}\par\smallskip

\noindent\begin{minipage}{\linewidth}\textbf{F1.088}
\begin{lstlisting}
1234122334 4143142213 1213122413 3144433514 2215123515 2435454535 2535225533 5555525555 5355235322
5344145342 2153215122 4122344413 3142213121 3122413414 4332214221 6123412233 4414314221 3151613661
3264463356 6236655666 5663666556 6556666666 6666666666 6666656666 6666666666 6666666666 6666666666
6666666666 6666666666 6666666666 6566666666 6666666666 6666556655 6663665666 5566326653 3644623166
3161513122 4134144332 2143216122 4122334414 3142213121 3122413314 4432214221 5123512243 5414435223
5325535555 5255555335 5225352535 4545342515 3215122415 3344413314 2213121312 2413414433 2214321
\end{lstlisting}
\end{minipage}\par\smallskip

\noindent\begin{minipage}{\linewidth}\textbf{F1.089}
\begin{lstlisting}
1234122334 4143142213 1213122413 3144433514 2215123515 2535454535 2535225532 5555335555 5355535322
5344145342 2153215122 4122344413 3142213121 3122413414 4332214321 6123412233 4414314221 6151613641
2264463356 6236655666 5653666553 6566636666 3666663666 6666656666 6666666666 6666666666 6666666666
6666666666 6666666666 6666666666 6566666666 3666663666 6366656355 6663565666 5566326653 3644622146
3161516122 4134144332 2143216123 4122334414 3142213121 3122413314 4432214221 5123512243 5414435223
5355535555 5335555235 5225352535 4545352515 3215122415 3344413314 2213121312 2413414433 2214321
\end{lstlisting}
\end{minipage}\par\smallskip

\noindent\begin{minipage}{\linewidth}\textbf{F1.090}
\begin{lstlisting}
1234122334 4143142213 1213122413 3144433514 2215123515 2535454535 2535225533 5455525555 5355235322
5344145342 2153215122 4122344413 3142213121 3122413414 4332214321 6123412233 4414314221 6151613661
3364463356 6256655666 5663666556 6566656666 6666666666 6666656666 6666666666 6666666666 6666666666
6666666666 6666666666 6666666666 6566666666 6666666666 6566656655 6663665666 5566526653 3644633166
3161516122 4134144332 2143216123 4122334414 3142213121 3122413314 4432214221 5123512243 5414435223
5325535555 5255545335 5225352535 4545352515 3215122415 3344413314 2213121312 2413414433 2214321
\end{lstlisting}
\end{minipage}\par\smallskip

\noindent\begin{minipage}{\linewidth}\textbf{F1.091}
\begin{lstlisting}
1234122334 4143142213 1213132413 3144432514 2315123515 2435454535 5535325533 5455535555 5355535322
5344145342 2153215132 4122344413 3142213121 3122413414 4332214321 6123412233 4414314221 6151613661
3364463356 6256655636 5663666553 6566656665 6666666666 6666656666 6666666666 6666666666 6666666666
6666666666 6666666666 6666666666 6566666666 6666666566 6566656355 6663665636 5566526653 3644633166
3161516122 4134144332 2143216123 4122334414 3142213121 3122413314 4432214231 5123512243 5414435223
5355535555 5355545335 5235355535 4545342515 3215132415 2344413314 2313121312 2413414433 2214321
\end{lstlisting}
\end{minipage}\par\smallskip

\noindent\begin{minipage}{\linewidth}\textbf{F1.092}
\begin{lstlisting}
1234122334 4143142213 1213132413 3144432514 2315123515 2535454535 5535325533 5455535545 5355535322
5344145342 2153215122 4122344413 3142213121 3122413414 4332214321 6123412233 4414316221 3151613661
3324463356 6256655636 5653666453 6566656665 6646666666 6636656666 6666666666 6666666666 6666666666
6666666666 6666666666 6666666666 6566366666 6666466566 6566656354 6663565636 5566526653 3644233166
3161513122 6134144332 2143216123 4122334414 3142213121 3122413314 4432214221 5123512243 5414435223
5355535545 5355545335 5235355535 4545352515 3215132415 2344413314 2313121312 2413414433 2214321
\end{lstlisting}
\end{minipage}\par\smallskip

\noindent\begin{minipage}{\linewidth}\textbf{F1.093}
\begin{lstlisting}
1234122334 4143142213 1213122413 3144433514 2215123515 2435454535 5535225533 5455335555 3355535322
5344145342 2153215132 4122344413 3142213121 3122413414 4332214221 6123412233 4414316221 3151613641
2266463356 6236655666 5643665456 6566636665 6666666666 6666656666 6666666666 6666666666 6666666666
6666666666 6666666666 6666666666 6566666666 6666666566 6366656654 5663465666 5566326653 3646622146
3161513122 6134144332 2143216122 4122334414 3142213121 3122413314 4432214231 5123512243 5414435223
5355533555 5335545335 5225355535 4545342515 3215122415 3344413314 2213121312 2413414433 2214321
\end{lstlisting}
\end{minipage}\par\smallskip

\noindent\begin{minipage}{\linewidth}\textbf{F1.094}
\begin{lstlisting}
1234122334 4143142213 1213122413 3144433514 2215123515 2535454535 2535225533 5455325545 5355235322
5344145342 2153215122 4122344413 3142213121 3122413414 4332214321 6123412233 4414314221 6151613661
3264463356 6256655666 5663666556 6566656666 6666666666 6666656666 6666666666 6666666666 6666666666
6666666666 6666666666 6666666666 6566666666 6666666666 6566656655 6663665666 5566526653 3644623166
3161516122 4134144332 2143216123 4122334414 3142213121 3122413314 4432214221 5123512243 5414435223
5325535545 5235545335 5225352535 4545352515 3215122415 3344413314 2213121312 2413414433 2214321
\end{lstlisting}
\end{minipage}\par\smallskip

\noindent\begin{minipage}{\linewidth}\textbf{F1.095}
\begin{lstlisting}
1234122334 4143142213 1213122413 3144433514 2315123515 2535454535 5535225533 5555325555 5355234322
5344145342 2153215122 4122344413 3142213121 3122413414 4332214321 6123412233 4414316221 3151613661
3364463356 6256655666 5663666556 6566666665 6666666666 6666656666 6666666666 6666666666 6666666666
6666666666 6666666666 6666666666 6566666666 6666666566 6666656655 6663665666 5566526653 3644633166
3161513122 6134144332 2143216123 4122334414 3142213121 3122413314 4432214221 5123512243 5414435223
4325535555 5235555335 5225355535 4545352515 3215132415 3344413314 2213121312 2413414433 2214321
\end{lstlisting}
\end{minipage}\par\smallskip

\noindent\begin{minipage}{\linewidth}\textbf{F1.096}
\begin{lstlisting}
1234122334 4143142213 1213122413 3144432514 2215123515 2535454535 5535325532 5455535545 3355235322
5344145342 2153215122 4122344413 3142213121 3122413414 4332214321 6123412233 4414314221 6151613641
3364463356 6236655636 5543665453 6556636665 3646663666 6336656666 6666666666 6666666666 6666666666
6666666666 6666666666 6666666666 6566336666 3666463566 6366556354 5663455636 5566326653 3644633146
3161516122 4134144332 2143216123 4122334414 3142213121 3122413314 4432214221 5123512243 5414435223
5325533545 5355545235 5235355535 4545352515 3215122415 2344413314 2213121312 2413414433 2214321
\end{lstlisting}
\end{minipage}\par\smallskip

\noindent\begin{minipage}{\linewidth}\textbf{F1.097}
\begin{lstlisting}
1234122334 4143142213 1213122413 3144432514 2215123515 2535454535 5535325532 5455535555 5355535322
5344145342 2153215122 4122344413 3142213121 3122413414 4332214221 6123412233 4414314221 6151613661
3364463356 6256655636 5653665556 6566656665 6646643666 6366656666 6666666666 6666666666 6666666666
6666666666 6666666666 6666666666 6566636666 3466466566 6566656655 5663565636 5566526653 3644633166
3161516122 4134144332 2143216122 4122334414 3142213121 3122413314 4432214221 5123512243 5414435223
5355535555 5355545235 5235355535 4545352515 3215122415 2344413314 2213121312 2413414433 2214321
\end{lstlisting}
\end{minipage}\par\smallskip

\noindent\begin{minipage}{\linewidth}\textbf{F1.098}
\begin{lstlisting}
1234122334 4143142213 1213122413 3144433514 2215123515 2435454535 2535225532 5555335555 5355535322
5344145342 2153215122 4122344413 3142213121 3122413414 4332214321 6123412233 4414314221 3151613661
2264463356 6256655666 5663666453 6566656666 6646666666 6366656666 6666666666 6666666666 6666666666
6666666666 6666666666 6666666666 6566636666 6666466666 6566656354 6663665666 5566526653 3644622166
3161513122 4134144332 2143216123 4122334414 3142213121 3122413314 4432214221 5123512243 5414435223
5355535555 5335555235 5225352535 4545342515 3215122415 3344413314 2213121312 2413414433 2214321
\end{lstlisting}
\end{minipage}\par\smallskip

\noindent\begin{minipage}{\linewidth}\textbf{F1.099}
\begin{lstlisting}
1234122334 4143142213 1213132413 3144432514 2315123515 2535454535 2535325532 5555335545 2355235322
5344145342 2153215122 4122344413 3142213121 3122413414 4332214221 6123412233 4414314221 6151613641
3364463356 6236655636 5543665453 6556636666 3646663666 6336656666 6666666666 6666666666 6666666666
6666666666 6666666666 6666666666 6566336666 3666463666 6366556354 5663455636 5566326653 3644633146
3161516122 4134144332 2143216122 4122334414 3142213121 3122413314 4432214221 5123512243 5414435223
5325532545 5335555235 5235352535 4545352515 3215132415 2344413314 2313121312 2413414433 2214321
\end{lstlisting}
\end{minipage}\par\smallskip

\noindent\begin{minipage}{\linewidth}\textbf{F1.100}
\begin{lstlisting}
1234122334 4143142213 1213122413 3144433514 2215123515 2535454535 5535225533 5455335545 5355535322
5344145342 2153215132 4122344413 3142213121 3122413414 4332214221 6123412233 4414314221 6151613641
3364463356 6236655636 5643665456 6556636665 3666666666 6336656666 6666666666 6666666666 6666666666
6666666666 6666666666 6666666666 6566336666 6666663566 6366556654 5663465636 5566326653 3644633146
3161516122 4134144332 2143216122 4122334414 3142213121 3122413314 4432214231 5123512243 5414435223
5355535545 5335545335 5225355535 4545352515 3215122415 3344413314 2213121312 2413414433 2214321
\end{lstlisting}
\end{minipage}\par\smallskip

\noindent\begin{minipage}{\linewidth}\textbf{F1.101}
\begin{lstlisting}
1234122334 4143142213 1213132413 3144432514 2215123515 2535454535 2535325533 5455235555 3355235322
5344145342 2153215122 4122344413 3142213121 3122413414 4332214321 6123412233 4414314221 3151613641
3326463356 6236655666 5643665453 6556636665 3646643666 6336656666 6666666666 6666666666 6666666666
6666666666 6666666666 6666666666 6566336666 3466463566 6366556354 5663465666 5566326653 3646233146
3161513122 4134144332 2143216123 4122334414 3142213121 3122413314 4432214221 5123512243 5414435223
5325533555 5325545335 5235352535 4545352515 3215122415 2344413314 2313121312 2413414433 2214321
\end{lstlisting}
\end{minipage}\par\smallskip

\noindent\begin{minipage}{\linewidth}\textbf{F1.102}
\begin{lstlisting}
1234122334 4143142213 1213122413 3144433514 2215123515 2535454535 5535225533 5555335545 5355535322
5344145342 2153215132 4122344413 3142213121 3122413414 4332214321 6123412233 4414316221 3151613661
2364463356 6256655636 5643666556 6566636665 3646663666 6336656666 6666666666 6666666666 6666666666
6666666666 6666666666 6666666666 6566336666 3666463566 6366656655 6663465636 5566526653 3644632166
3161513122 6134144332 2143216123 4122334414 3142213121 3122413314 4432214231 5123512243 5414435223
5355535545 5335555335 5225355535 4545352515 3215122415 3344413314 2213121312 2413414433 2214321
\end{lstlisting}
\end{minipage}\par\smallskip

\noindent\begin{minipage}{\linewidth}\textbf{F1.103}
\begin{lstlisting}
1234122334 4143142213 1213122413 3144432514 2315123515 2535454535 5535225532 5555335555 5355235322
5344145342 2153215132 4122344413 3142213121 3122413414 4332214321 6123412233 4414314221 3151613641
3366463356 6236655636 5543665453 6556636665 3646643666 6336656666 6666666666 6666666666 6666666666
6666666666 6666666666 6666666666 6566336666 3466463566 6366556354 5663455636 5566326653 3646633146
3161513122 4134144332 2143216123 4122334414 3142213121 3122413314 4432214231 5123512243 5414435223
5325535555 5335555235 5225355535 4545352515 3215132415 2344413314 2213121312 2413414433 2214321
\end{lstlisting}
\end{minipage}\par\smallskip

\noindent\begin{minipage}{\linewidth}\textbf{F1.104}
\begin{lstlisting}
1234122334 4143142213 1213122413 3144433514 2215123515 2535454535 2535325533 5455335555 3355535322
5344145342 2153215132 4122344413 3142213121 3122413414 4332214321 6123412233 4414314221 3151613641
2364463356 6236655636 5663665553 6566636666 3646643666 6636656666 6666666666 6666666666 6666666666
6666666666 6666666666 6666666666 6566366666 3466463666 6366656355 5663665636 5566326653 3644632146
3161513122 4134144332 2143216123 4122334414 3142213121 3122413314 4432214231 5123512243 5414435223
5355533555 5335545335 5235352535 4545352515 3215122415 3344413314 2213121312 2413414433 2214321
\end{lstlisting}
\end{minipage}\par\smallskip

\noindent\begin{minipage}{\linewidth}\textbf{F1.105}
\begin{lstlisting}
1234122334 4143142213 1213122413 3144433514 2215123515 2535454535 5535325533 5555335555 5355534322
5344145342 2153215132 4122344413 3142213121 3122413414 4332214321 6123412233 4414314221 6151613661
2264463356 6236655636 5663666456 6566536666 3666643666 6366656666 6666666666 6666666666 6666666666
6666666666 6666666666 6666666666 6566636666 3466663666 6356656654 6663665636 5566326653 3644622166
3161516122 4134144332 2143216123 4122334414 3142213121 3122413314 4432214231 5123512243 5414435223
4355535555 5335555335 5235355535 4545352515 3215122415 3344413314 2213121312 2413414433 2214321
\end{lstlisting}
\end{minipage}\par\smallskip

\noindent\begin{minipage}{\linewidth}\textbf{F1.106}
\begin{lstlisting}
1234122334 4143142213 1213122413 3144433514 2315123515 2435454535 5535325532 5455335555 3355535322
5344145342 2153215132 4122344413 3142213121 3122413414 4332214221 6123412233 4414314221 3151613641
2266463356 6236655636 5543665453 6556656665 3646646666 6336656666 6666666666 6666666666 6666666666
6666666666 6666666666 6666666666 6566336666 6466463566 6566556354 5663455636 5566326653 3646622146
3161513122 4134144332 2143216122 4122334414 3142213121 3122413314 4432214231 5123512243 5414435223
5355533555 5335545235 5235355535 4545342515 3215132415 3344413314 2213121312 2413414433 2214321
\end{lstlisting}
\end{minipage}\par\smallskip

\noindent\begin{minipage}{\linewidth}\textbf{F1.107}
\begin{lstlisting}
1234122334 4143142213 1213122413 3144432514 2215123515 2435454535 2535225532 5455335555 2355235322
5344145342 2153215122 4122344413 3142213121 3122413414 4332214221 6123412233 4414316221 3151613651
3366463356 6236655636 5553665553 6556636665 3644663666 6336656666 6666666666 6666666666 6666666666
6666666666 6666666666 6666666666 6566336666 3664463566 6366556355 5663555636 5566326653 3646633156
3161513122 6134144332 2143216122 4122334414 3142213121 3122413314 4432214221 5123512243 5414435223
5325532555 5335545235 5225352535 4545342515 3215122415 2344413314 2213121312 2413414433 2214321
\end{lstlisting}
\end{minipage}\par\smallskip

\subsection*{Family $\mathcal F_2$ (14 partitions)}

\noindent\begin{minipage}{\linewidth}\textbf{F2.001}
\begin{lstlisting}
1234122334 4143142213 1213122413 3144433514 2215123515 2435454535 5535225532 5555335555 5355535322
5344145342 2153215122 4122344413 3142213121 3122413414 4332214321 6123412233 4414314221 3151613641
2266463356 6236655666 5643665656 6566666666 6666666666 6666156666 6666666666 6666666666 6666666666
6666666666 6666666666 6666666666 6516666666 6666666666 6666656656 5663465666 5566326653 3646622146
3161513122 4134144332 2143216123 4122334414 3142213121 3122413314 4432214221 5123512243 5414435223
5355535555 5335555235 5225355535 4545342515 3215122415 3344413314 2213121312 2413414433 2214321
\end{lstlisting}
\end{minipage}\par\smallskip

\noindent\begin{minipage}{\linewidth}\textbf{F2.002}
\begin{lstlisting}
1234122334 4143142213 1213122413 3144433514 2215123515 2535454535 5535325533 5555335555 5355535322
5344145342 2153215122 4122344413 3142213121 3122413414 4332214321 6123412233 4414314221 6151613661
3264463356 6256655666 5663666556 6566636666 6666666666 6666156666 6666666666 6666666666 6666666666
6666666666 6666666666 6666666666 6516666666 6666666666 6366656655 6663665666 5566526653 3644623166
3161516122 4134144332 2143216123 4122334414 3142213121 3122413314 4432214221 5123512243 5414435223
5355535555 5335555335 5235355535 4545352515 3215122415 3344413314 2213121312 2413414433 2214321
\end{lstlisting}
\end{minipage}\par\smallskip

\noindent\begin{minipage}{\linewidth}\textbf{F2.003}
\begin{lstlisting}
1234122334 4143142213 1213122413 3144432514 2215123515 2535454535 5535225533 5555535555 5355535322
5344145342 2153215122 4122344413 3142213121 3122413414 4332214221 6123412233 4414314221 3151613641
3364463356 6256655666 5663666656 6566666666 3646643666 6336156666 6666666666 6666666666 6666666666
6666666666 6666666666 6666666666 6516336666 3466463666 6666656656 6663665666 5566526653 3644633146
3161513122 4134144332 2143216122 4122334414 3142213121 3122413314 4432214221 5123512243 5414435223
5355535555 5355555335 5225355535 4545352515 3215122415 2344413314 2213121312 2413414433 2214321
\end{lstlisting}
\end{minipage}\par\smallskip

\noindent\begin{minipage}{\linewidth}\textbf{F2.004}
\begin{lstlisting}
1234122334 4143142213 1213132413 3144432514 2315123515 2535454535 5535325532 5555535555 3355535322
5344145342 2153215122 4122344413 3142213121 3122413414 4332214321 6123412233 4414314221 6151613641
3364463356 6236655636 5553665453 6556636666 6646666666 6336156666 6666666666 6666666666 6666666666
6666666666 6666666666 6666666666 6516336666 6666466666 6366556354 5663555636 5566326653 3644633146
3161516122 4134144332 2143216123 4122334414 3142213121 3122413314 4432214221 5123512243 5414435223
5355533555 5355555235 5235355535 4545352515 3215132415 2344413314 2313121312 2413414433 2214321
\end{lstlisting}
\end{minipage}\par\smallskip

\noindent\begin{minipage}{\linewidth}\textbf{F2.005}
\begin{lstlisting}
1234122334 4143142213 1213122413 3144433514 2215123515 2535454535 5535225533 5455325545 5355535322
5344145342 2153215122 4122344413 3142213121 3122413414 4332214221 6123412233 4414314221 3151613641
3364463356 6236655666 5643665556 6556656666 3646663666 6636156666 6666666666 6666666666 6666666666
6666666666 6666666666 6666666666 6516366666 3666463666 6566556655 5663465666 5566326653 3644633146
3161513122 4134144332 2143216122 4122334414 3142213121 3122413314 4432214221 5123512243 5414435223
5355535545 5235545335 5225355535 4545352515 3215122415 3344413314 2213121312 2413414433 2214321
\end{lstlisting}
\end{minipage}\par\smallskip

\noindent\begin{minipage}{\linewidth}\textbf{F2.006}
\begin{lstlisting}
1234122334 4143142213 1213122413 3144433514 2215123515 2435454535 5535225532 5555335545 5355535322
5344145342 2153215122 4122344413 3142213121 3122413414 4332214221 6123412233 4414314221 6151613641
2264463356 6236655666 5543665453 6556636665 6666666666 6636156666 6666666666 6666666666 6666666666
6666666666 6666666666 6666666666 6516366666 6666666566 6366556354 5663455666 5566326653 3644622146
3161516122 4134144332 2143216122 4122334414 3142213121 3122413314 4432214221 5123512243 5414435223
5355535545 5335555235 5225355535 4545342515 3215122415 3344413314 2213121312 2413414433 2214321
\end{lstlisting}
\end{minipage}\par\smallskip

\noindent\begin{minipage}{\linewidth}\textbf{F2.007}
\begin{lstlisting}
1234122334 4143142213 1213122413 3144433514 2215123515 2435454535 5535225533 5555525555 5355535322
5344145342 2153215122 4122344413 3142213121 3122413414 4332214321 6123412233 4414314221 6151613641
3364463356 6236655666 5643665453 6566656665 3646663666 6636156666 6666666666 6666666666 6666666666
6666666666 6666666666 6666666666 6516366666 3666463566 6566656354 5663465666 5566326653 3644633146
3161516122 4134144332 2143216123 4122334414 3142213121 3122413314 4432214221 5123512243 5414435223
5355535555 5255555335 5225355535 4545342515 3215122415 3344413314 2213121312 2413414433 2214321
\end{lstlisting}
\end{minipage}\par\smallskip

\noindent\begin{minipage}{\linewidth}\textbf{F2.008}
\begin{lstlisting}
1234122334 4143142213 1213122413 3144433514 2315123515 2435454535 5535325532 5455535555 3355535322
5344145342 2153215132 4122344413 3142213121 3122413414 4332214221 6123412233 4414314221 6151613641
3264463356 6236655666 5563665653 6566636666 3646646666 6666156666 6666666666 6666666666 6666666666
6666666666 6666666666 6666666666 6516666666 6466463666 6366656356 5663655666 5566326653 3644623146
3161516122 4134144332 2143216122 4122334414 3142213121 3122413314 4432214231 5123512243 5414435223
5355533555 5355545235 5235355535 4545342515 3215132415 3344413314 2213121312 2413414433 2214321
\end{lstlisting}
\end{minipage}\par\smallskip

\noindent\begin{minipage}{\linewidth}\textbf{F2.009}
\begin{lstlisting}
1234122334 4143142213 1213122413 3144433514 2215123515 2535454535 5535225533 5455525555 3355535322
5344145342 2153215132 4122344413 3142213121 3122413414 4332214321 6123412233 4414316221 3151613651
3266463356 6236655666 5653666653 6566636665 3664643666 6636156666 6666666666 6666666666 6666666666
6666666666 6666666666 6666666666 6516366666 3464663566 6366656356 6663565666 5566326653 3646623156
3161513122 6134144332 2143216123 4122334414 3142213121 3122413314 4432214231 5123512243 5414435223
5355533555 5255545335 5225355535 4545352515 3215122415 3344413314 2213121312 2413414433 2214321
\end{lstlisting}
\end{minipage}\par\smallskip

\noindent\begin{minipage}{\linewidth}\textbf{F2.010}
\begin{lstlisting}
1234122334 4143142213 1213132413 3144432514 2215123515 2535454535 5535225533 5555335555 5355535322
5344145342 2153215122 4122344413 3142213121 3122413414 4332214321 6123412233 4414316221 3151613661
3324463356 6256655666 5553665553 6556656665 6666643666 6366156666 6666666666 6666666666 6666666666
6666666666 6666666666 6666666666 6516636666 3466666566 6566556355 5663555666 5566526653 3644233166
3161513122 6134144332 2143216123 4122334414 3142213121 3122413314 4432214221 5123512243 5414435223
5355535555 5335555335 5225355535 4545352515 3215122415 2344413314 2313121312 2413414433 2214321
\end{lstlisting}
\end{minipage}\par\smallskip

\noindent\begin{minipage}{\linewidth}\textbf{F2.011}
\begin{lstlisting}
1234122334 4143142213 1213122413 3144433514 2215123515 2435454535 5535225533 5555325555 3355235322
5344145342 2153215122 4122344413 3142213121 3122413414 4332214221 6123412233 4414314221 3151613641
3264463356 6236655636 5543665453 6566656666 3646663666 6666156666 6666666666 6666666666 6666666666
6666666666 6666666666 6666666666 6516666666 3666463666 6566656354 5663455636 5566326653 3644623146
3161513122 4134144332 2143216122 4122334414 3142213121 3122413314 4432214221 5123512243 5414435223
5325533555 5235555335 5225355535 4545342515 3215122415 3344413314 2213121312 2413414433 2214321
\end{lstlisting}
\end{minipage}\par\smallskip

\noindent\begin{minipage}{\linewidth}\textbf{F2.012}
\begin{lstlisting}
1234122334 4143142213 1213122413 3144433514 2315123515 2535454535 2535325533 5555535545 3355535322
5344145342 2153215122 4122344413 3142213121 3122413414 4332214321 6123412233 4414314221 3151613661
2264463356 6256655666 5553666556 6556656665 3646663666 6336156666 6666666666 6666666666 6666666666
6666666666 6666666666 6666666666 6516336666 3666463566 6566556655 6663555666 5566526653 3644622166
3161513122 4134144332 2143216123 4122334414 3142213121 3122413314 4432214221 5123512243 5414435223
5355533545 5355555335 5235352535 4545352515 3215132415 3344413314 2213121312 2413414433 2214321
\end{lstlisting}
\end{minipage}\par\smallskip

\noindent\begin{minipage}{\linewidth}\textbf{F2.013}
\begin{lstlisting}
1234122334 4143142213 1213132413 3144432514 2215123515 2535454535 5535225532 5455535555 5355535322
5344145342 2153215132 4122344413 3142213121 3122413414 4332214221 6123412233 4414316221 3151613651
3366463356 6256655666 5553665553 6556656665 6644663666 6336156666 6666666666 6666666666 6666666666
6666666666 6666666666 6666666666 6516336666 3664466566 6566556355 5663555666 5566526653 3646633156
3161513122 6134144332 2143216122 4122334414 3142213121 3122413314 4432214231 5123512243 5414435223
5355535555 5355545235 5225355535 4545352515 3215122415 2344413314 2313121312 2413414433 2214321
\end{lstlisting}
\end{minipage}\par\smallskip

\noindent\begin{minipage}{\linewidth}\textbf{F2.014}
\begin{lstlisting}
1234122334 4143142213 1213122413 3144432514 2215123515 2535454535 2535225532 5455335545 2355234322
5344145342 2153215122 4122344413 3142213121 3122413414 4332214221 6123412233 4414316221 3151613651
3366463356 6236655636 5553665553 6556636665 3644663666 6336156666 6666666666 6666666666 6666666666
6666666666 6666666666 6666666666 6516336666 3664463566 6366556355 5663555636 5566326653 3646633156
3161513122 6134144332 2143216122 4122334414 3142213121 3122413314 4432214221 5123512243 5414435223
4325532545 5335545235 5225352535 4545352515 3215122415 2344413314 2213121312 2413414433 2214321
\end{lstlisting}
\end{minipage}\par\smallskip

\subsection*{Family $\mathcal F_3$ (5 partitions)}

\noindent\begin{minipage}{\linewidth}\textbf{F3.001}
\begin{lstlisting}
1234122334 4143142213 1213122413 3144432514 2215123515 2535454535 5535225532 5555535555 5355535322
5344145342 2153215122 4122344413 3142213121 3122413414 4332214321 6123412233 4414316221 3151613641
3366463356 6256655666 5643666656 6566666665 6666666666 6666656666 6666666666 6666666666 6666666666
6666666666 6666666666 6666666666 6566666666 6666666566 6666656656 6663565666 5566526653 3646633146
3161513122 6134144332 2143216123 4122334414 3142213121 3122413314 4432214231 5123512243 5414435223
5355535555 5355555235 5235355535 4545352515 3215132415 2344413314 2213121312 2413414433 2214321
\end{lstlisting}
\end{minipage}\par\smallskip

\noindent\begin{minipage}{\linewidth}\textbf{F3.002}
\begin{lstlisting}
1234122334 4143142213 1213122413 3144433514 2215123515 2535454535 5535225532 5555335545 5355535322
5344145342 2153215122 4122344413 3142213121 3122413414 4332214221 6123412233 4414314221 3151613641
3366463356 6256655636 5553666656 6556666666 6666666666 6666656666 6666666666 6666666666 6666666666
6666666666 6666666666 6666666666 6566666666 6666666666 6666556656 6663555636 5566526653 3646633146
3161513122 4134144332 2143216123 4122334414 3142213121 3122413314 4432214221 5123512243 5414435223
5355535545 5335555235 5225355535 4545342515 3215122415 3344413314 2213121312 2413414433 2214321
\end{lstlisting}
\end{minipage}\par\smallskip

\noindent\begin{minipage}{\linewidth}\textbf{F3.003}
\begin{lstlisting}
1234122334 4143142213 1213132413 3144432514 2215123515 2435454535 5535325532 5455335545 5355535322
5344145342 2153215122 4122344413 3142213121 3122413414 4332214321 6123412233 4414316221 3151613651
3366463356 6256655666 5653666553 6566666665 6664666666 6666656666 6666666666 6666666666 6666666666
6666666666 6666666666 6666666666 6566666666 6664666566 6566656655 6663555666 5566326653 3646633156
3161513122 6134144332 2143216123 4122334414 3142213121 3122413314 4432214221 5123512243 5414435223
5355535545 5335545235 5235355535 4545342515 3215132415 2344413314 2313121312 2413414433 2214321
\end{lstlisting}
\end{minipage}\par\smallskip

\noindent\begin{minipage}{\linewidth}\textbf{F3.004}
\begin{lstlisting}
1234122334 4143142213 1213122413 3144433514 2215123515 2535454535 5535225533 5455535555 5355535322
5344145342 2153215122 4122344413 3142213121 3122413414 4332214221 6123412233 4414314221 6151613651
2364463356 6236655666 5663665656 6566656666 6666666666 6666656666 6666666666 6666666666 6666666666
6666666666 6666666666 6666666666 6566666666 6666666666 6566656656 5663655636 5566326653 3644622156
3161516122 4134144332 2143216123 4122334414 3142213121 3122413314 4432214231 5123512243 5414435223
5355535555 5355545335 5225355535 4545352515 3215122415 3344413314 2213121312 2413414433 2214321
\end{lstlisting}
\end{minipage}\par\smallskip

\noindent\begin{minipage}{\linewidth}\textbf{F3.005}
\begin{lstlisting}
1234122334 4143142213 1213122413 3144433514 2215123515 2535454535 5535225533 5455525555 5355535322
5344145342 2153215122 4122344413 3142213121 3122413414 4332214221 6123412233 4414314221 6151613641
3364463356 6236655636 5653665453 6556656666 3646666666 6666656666 6666666666 6666666666 6666666666
6666666666 6666666666 6666666666 6566666666 6666466666 6566656654 6663465636 5566526653 3644633166
3161516122 4134144332 2143216122 4122334414 3142213121 3122413314 4432214231 5123512243 5414435223
5355535555 5235555335 5225352535 4545352515 3215132415 3344413314 2213121312 2413414433 2214321
\end{lstlisting}
\end{minipage}\par\smallskip

\subsection*{Family $\mathcal F_4$ (4 partitions)}

\noindent\begin{minipage}{\linewidth}\textbf{F4.001}
\begin{lstlisting}
1234122334 4143142213 1213122413 3144432514 2215123515 2435454535 2535225532 5555535545 2355235322
5344145342 2153215132 4122344413 3142213121 3122413414 4332214321 6123412233 4414316221 3151613641
3364463356 6236655666 5643666456 6566636665 6666666666 6666156666 6666666666 6666666666 6666666646
6666666666 6666666666 6666666666 6516666666 6666666566 6366656654 6663465666 5566326653 3644633146
3161513122 6134144332 2143216123 4122334414 3142213121 3122413314 4432214231 5123512243 5414435223
5325532545 5355555235 5225352535 4545342515 3215122415 2344413314 2213121312 2413414433 2214321
\end{lstlisting}
\end{minipage}\par\smallskip

\noindent\begin{minipage}{\linewidth}\textbf{F4.002}
\begin{lstlisting}
1234122334 4143142213 1213122413 3144432514 2315123515 2535454535 5535325533 5455535545 2355235322
5344145342 2153215132 4122344413 3142213121 3122413414 4332214221 6123412233 4414314221 3151613641
3364463356 6236655636 5643665556 6566666666 3666666666 6666156336 6636666666 6666666666 6666666666
6666666666 6666666666 6666366633 6516666666 6666663666 6666656655 5663465636 5566326653 3644633146
3161513122 4134144332 2143216122 4122334414 3142213121 3122413314 4432214231 5123512243 5414435223
5325532545 5355545335 5235355535 4545352515 3215132415 2344413314 2213121312 2413414433 2214321
\end{lstlisting}
\end{minipage}\par\smallskip

\noindent\begin{minipage}{\linewidth}\textbf{F4.003}
\begin{lstlisting}
1234122334 4143142213 1213122413 3144432514 2315123515 2535454535 5535225533 5455335555 2355235322
5344145342 2153215132 4122344413 3142213121 3122413414 4332214321 6123412233 4414314221 3151613641
3366463356 6236655666 5643666653 6566636666 6646666666 6666151636 6666666666 6666661616 6666666646
6666666616 1666666666 6666666636 1516666666 6666466666 6366656356 6663465666 5566326653 3646633146
3161513122 4134144332 2143216123 4122334414 3142213121 3122413314 4432214231 5123512243 5414435223
5325532555 5335545335 5225355535 4545352515 3215132415 2344413314 2213121312 2413414433 2214321
\end{lstlisting}
\end{minipage}\par\smallskip

\noindent\begin{minipage}{\linewidth}\textbf{F4.004}
\begin{lstlisting}
1234122334 4143142213 1213122413 3144433514 2215123515 2435454535 5535225532 5555335555 3355235322
5344145342 2153215122 4122344413 3142213121 3122413414 4332214321 6123412233 4414314221 3151613641
2264463356 6236655636 5543665453 6556636665 3646643666 6336651336 6663664636 6666361613 3666636646
6366663316 1636666636 4663666633 1566336666 3466463566 6366556354 5663455636 5566326653 3644622146
3161513122 4134144332 2143216123 4122334414 3142213121 3122413314 4432214221 5123512243 5414435223
5325533555 5335555235 5225355535 4545342515 3215122415 3344413314 2213121312 2413414433 2214321
\end{lstlisting}
\end{minipage}\par\smallskip

\subsection*{Family $\mathcal F_5$ (1 partitions)}

\noindent\begin{minipage}{\linewidth}\textbf{F5.001}
\begin{lstlisting}
1234415555 5553164332 6142213632 1555215143 4313126441 3555552614 4213632361 6552555231 3122441223
6325555555 3634461662 2463655555 6541223634 3221612555 5551662446 3632443525 5553634362 1614263132
3555546363 2343622144 5655355161 4263131413 4155553534 3662144163 2143551555 3131223412 3614455555
5544163214 3221313555 5553412361 4412663435 3555514314 1313624161 5535565441 2263442363 6455553231
3624161263 4363551525 3442363644 2661535555 2161223436 3221456555 5563642266 1644363555 5555236322
1442213132 5552556163 2363124416 2555553144 6213134341 5125555236 3122416233 4613555555 5144321
\end{lstlisting}
\end{minipage}\par\smallskip

\end{document}
